\documentclass{article}
\usepackage[margin=0.8in]{geometry}
\usepackage{amsmath,amssymb,amsthm,mathtools}
\usepackage{mathrsfs,bm,bbm}
\usepackage{graphicx,booktabs,array,multirow,tabularx,longtable}
\usepackage{enumitem}
\usepackage{xcolor}
\usepackage{algorithm,algpseudocode}
\usepackage{subcaption,tikz}
\usepackage{siunitx,cancel,accents}
\usepackage{microtype}
\usepackage[numbers,sort&compress]{natbib}
\usepackage[colorlinks=true,linkcolor=blue,citecolor=blue,urlcolor=blue]{hyperref}
\usepackage{cleveref}
\setlength{\emergencystretch}{2em}
\newtheorem{assumption}{Assumption}
\newtheorem{definition}{Definition}
\newtheorem{theorem}{Theorem}
\newtheorem{lemma}{Lemma}
\newtheorem{proposition}{Proposition}
\newtheorem{openendedquestion}{Open-ended Question}
\newtheorem{conjecture}{Conjecture}
\newcommand{\coreref}[1]{\ref{#1}}

\begin{document}

\sloppy

\section*{eid\_\allowbreak{}proxychain\_\allowbreak{}cubic\_\allowbreak{}inverse}

\noindent\textit{Specialization:} v1\par

\paragraph{TL;DR.} A positive-skew, one-global-confounder three-node chain is globally invertible from its third cumulant tensor alone. The unique rank-two slice locates the root, its signed kernel locates the middle node, and two adjugate contractions recover both edges with uniform denominator margins. On the exact published proxy graph, an independent full-ideal replication and re-audit tests the additional localized relation \(\beta W_2^{\mathrm{el}}-W_3^{\mathrm{el}}=0\), with zero remainder against the pinned entire basis as the first checkpoint; a complete real exceptional-fiber atlas remains an explicit open question. An opposite-chain covariance overlap proves that skewness supplies genuinely new order information, while simultaneous bounded-moment inversion gives finite-sample uniform coverage and root-\(n\) diameter without nonsingular influence covariance.

\section{Project justification}

\paragraph{Gap.} The closest proxy-LvLiNGAM result uses fourth-order cumulants for identification and reports nonidentification through order three on the exact one-proxy graph. The status of the independently derived localized relation \(\beta W_2^{\mathrm{el}}-W_3^{\mathrm{el}}=0\) under the released full elimination basis is not reported, and existing cumulant rank methods neither resolve the compact positive chamber nor classify real failures of rational charts.

\paragraph{Niche.} The joint three-variable cubic tensor has chain-specific rank and sign geometry that is invisible to bivariate root enumeration and can collapse both label and coefficient ambiguity.

\paragraph{Fill.} The project derives a constant-cost adjugate recovery map on the positive compact chamber with global denominator margins, independently replicates and re-audits the exact published full ideal with a pinned zero-remainder gate, proves a covariance-only open overlap, attaches honest finite-sample inversion, and isolates complete real-fiber stratification as the first-class unresolved theorem question.

\section{Related work}

Tramontano et al. (2025), Theorems 3.5--3.6, Figure 3, Appendix C.1.1, and the released \(\texttt{NonGaussianIdentifiability.m2}\) file provide the exact one-proxy graph, the fourth-order estimator, and the reported third-order boundary. The localized theorem is an independent full-ideal replication and re-audit of that boundary: it tests the specifically alleged additional relation \(\beta W_2^{\mathrm{el}}-W_3^{\mathrm{el}}=0\) and does not attribute a restricted proof strategy to the published argument. Cai et al. (2023), Theorem 2 and Appendix A.1, motivate closed-form cumulant extraction for a one-latent component but use covariance with a selected higher order. Schkoda, Robeva, and Drton (2024), their source-rank theorem, motivate the rank diagnostic but stack orders three and four for one latent parent. Garcia-Puente, Spielvogel, and Sullivant (2010), Propositions 2--4, supply the full-elimination-ideal criterion used in the independent re-audit. Drton, Foygel, and Sullivant (2011) supply the global-versus-generic fiber distinction used to formulate the exceptional atlas. Ghassami et al. (2020) motivate comparison of entire model images, while Foygel Barber et al. (2022) treat rational covariance recovery for known explicit-latent graphs rather than cross-label overlap. Babii (2020) and Andrews and Guggenberger (2019) inform the confidence-set inversion, but neither treats exact bounded cubic charts with discrete causal-order selection.

\section{Setup}

Target (estimand / structural parameter): \(B=B_\pi(u,v)\) on the unknown-order principal real cubic model; the compact probabilistic delivery satisfies \(\{C_3(P):P\in\mathcal P_+\}\subset\mathcal C^\circ\), while no equality is asserted on the unrestricted published class..
Primary functional / estimator / diagnostic: \(\operatorname{Rec}(C)=B_\pi(u,v)\) for \(C\in\mathcal C^\circ\); in particular, \(\operatorname{Rec}(C_3(P))=B_\pi(u,v)\) for every \(P\in\mathcal P_+\)..
\begin{itemize}
  \item \texttt{\textbackslash{}\allowbreak{}mathcal V} --- \texttt{finite label set}
  \par\smallskip\noindent\emph{Definition.} \(\mathcal V=\{1,2,3\}\)
  \item \texttt{I\_\allowbreak{}3} --- \texttt{matrix}; \(\mathbb R^{3\times 3}\)
  \par\smallskip\noindent\emph{Definition.} \texttt{the identity matrix}
  \item \texttt{e\_\allowbreak{}i} --- vector indexed by \(i\in\mathcal V\); \(\mathbb R^3\)
  \par\smallskip\noindent\emph{Definition.} the \(i\)-th coordinate vector
  \item \texttt{\textbackslash{}\allowbreak{}mathfrak S\_\allowbreak{}3} --- \texttt{finite set}
  \par\smallskip\noindent\emph{Definition.} the permutations of \(\mathcal V\)
  \item \texttt{\textbackslash{}\allowbreak{}pi} --- \texttt{permutation}; \(\mathfrak S_3\)
  \item \texttt{\textbackslash{}\allowbreak{}pi\_\allowbreak{}f} --- \texttt{permutation}; \(\mathfrak S_3\)
  \par\smallskip\noindent\emph{Definition.} \(\pi_f=(1,2,3)\)
  \item \texttt{\textbackslash{}\allowbreak{}pi\_\allowbreak{}r} --- \texttt{permutation}; \(\mathfrak S_3\)
  \par\smallskip\noindent\emph{Definition.} \(\pi_r=(3,2,1)\)
  \item \texttt{a} --- \texttt{label}; \(\mathcal V\)
  \par\smallskip\noindent\emph{Definition.} the root coordinate in \(\pi=(a,b,c)\)
  \item \texttt{b} --- \texttt{label}; \(\mathcal V\)
  \par\smallskip\noindent\emph{Definition.} the middle coordinate in \(\pi=(a,b,c)\)
  \item \texttt{c} --- \texttt{label}; \(\mathcal V\)
  \par\smallskip\noindent\emph{Definition.} the terminal coordinate in \(\pi=(a,b,c)\)
  \item \texttt{u} --- \texttt{scalar}; \(\mathbb R\); direct effect from \(a\) to \(b\)
  \item \texttt{v} --- \texttt{scalar}; \(\mathbb R\); direct effect from \(b\) to \(c\)
  \item \texttt{\textbackslash{}\allowbreak{}gamma} --- \texttt{vector}; \(\mathbb R^3\); \texttt{global latent loading}
  \item \texttt{\textbackslash{}\allowbreak{}gamma\textasciicircum{}\{\textbackslash{}\allowbreak{}mathrm\{rev\}\allowbreak{}\}\allowbreak{}} --- \texttt{vector}; \(\mathbb R^3\)
  \par\smallskip\noindent\emph{Definition.} the coordinate reversal of \(\gamma\)
  \item \texttt{s} --- \texttt{vector}; \((0,\infty)^3\); \texttt{idiosyncratic residual scales}
  \item \texttt{s\textasciicircum{}\{\textbackslash{}\allowbreak{}mathrm\{rev\}\allowbreak{}\}\allowbreak{}} --- \texttt{vector}; \((0,\infty)^3\)
  \par\smallskip\noindent\emph{Definition.} the coordinate reversal of \(s\)
  \item \texttt{H} --- \texttt{random variable}; \(\mathbb R\); \texttt{global latent source}
  \item \texttt{\textbackslash{}\allowbreak{}xi} --- \texttt{random vector}; \(\mathbb R^3\); \texttt{idiosyncratic sources}
  \item \texttt{\textbackslash{}\allowbreak{}kappa} --- \texttt{vector}; \(\mathbb R^4\)
  \par\smallskip\noindent\emph{Definition.} \(\kappa=(\operatorname{Cum}_3(H),\operatorname{Cum}_3(\xi_1),\operatorname{Cum}_3(\xi_2),\operatorname{Cum}_3(\xi_3))\)
  \item \texttt{\textbackslash{}\allowbreak{}kappa\_\allowbreak{}0} --- \texttt{scalar}; \([1/2,1]\)
  \item \texttt{b\_\allowbreak{}0} --- \texttt{positive scalar}
  \par\smallskip\noindent\emph{Definition.} \(b_0=(\kappa_0+\sqrt{\kappa_0^2+4})/2\)
  \item \texttt{a\_\allowbreak{}0} --- \texttt{positive scalar}
  \par\smallskip\noindent\emph{Definition.} \(a_0=1/b_0\)
  \item \texttt{B\_\allowbreak{}\textbackslash{}\allowbreak{}pi} --- \texttt{matrix}; \(\mathbb R^{3\times3}\)
  \par\smallskip\noindent\emph{Definition.} \(B_\pi(u,v)=u e_b e_a^{\mathsf T}+v e_c e_b^{\mathsf T}\)
  \item \texttt{A\_\allowbreak{}\textbackslash{}\allowbreak{}pi} --- \texttt{matrix}; \(\mathbb R^{3\times3}\)
  \par\smallskip\noindent\emph{Definition.} \(A_\pi=(I_3-B_\pi)^{-1}\)
  \item \texttt{X} --- \texttt{random vector}; \(\mathbb R^3\); \texttt{observed vector}
  \item \texttt{P} --- \texttt{probability law}; laws on \(\mathbb R^3\)
  \par\smallskip\noindent\emph{Definition.} \(P=\mathcal L(X)\)
  \item \texttt{C} --- \texttt{symmetric tensor}; \(\operatorname{Sym}^3(\mathbb R^3)\)
  \par\smallskip\noindent\emph{Definition.} \(C=C_3(P)=\operatorname{Cum}_3(X)\)
  \item \texttt{M\_\allowbreak{}i} --- matrix indexed by \(i\in\mathcal V\); \(\mathbb R^{3\times3}\)
  \par\smallskip\noindent\emph{Definition.} \(M_i=C(e_i,\cdot,\cdot)\)
  \item \texttt{x} --- \texttt{scalar}
  \par\smallskip\noindent\emph{Definition.} \(x=\gamma_a\)
  \item \texttt{y} --- \texttt{scalar}
  \par\smallskip\noindent\emph{Definition.} \(y=\gamma_b\)
  \item \texttt{z} --- \texttt{scalar}
  \par\smallskip\noindent\emph{Definition.} \(z=\gamma_c\)
  \item \texttt{p} --- \texttt{scalar}
  \par\smallskip\noindent\emph{Definition.} \(p=ux+y\)
  \item \texttt{q} --- \texttt{scalar}
  \par\smallskip\noindent\emph{Definition.} \(q=vp+z\)
  \item \texttt{h} --- \texttt{scalar}
  \par\smallskip\noindent\emph{Definition.} \(h=\operatorname{Cum}_3(H)\)
  \item \texttt{r} --- \texttt{scalar}
  \par\smallskip\noindent\emph{Definition.} \(r=s_a^3\operatorname{Cum}_3(\xi_a)\)
  \item \texttt{t} --- \texttt{scalar}
  \par\smallskip\noindent\emph{Definition.} \(t=s_b^3\operatorname{Cum}_3(\xi_b)\)
  \item \texttt{w} --- \texttt{scalar}
  \par\smallskip\noindent\emph{Definition.} \(w=s_c^3\operatorname{Cum}_3(\xi_c)\)
  \item \texttt{g} --- \texttt{vector}; \(\mathbb R^3\)
  \par\smallskip\noindent\emph{Definition.} \(g=(x,p,q)^{\mathsf T}\)
  \item \texttt{d} --- \texttt{vector}; \(\mathbb R^3\)
  \par\smallskip\noindent\emph{Definition.} \(d=(1,u,uv)^{\mathsf T}\)
  \item \texttt{e} --- \texttt{vector}; \(\mathbb R^3\)
  \par\smallskip\noindent\emph{Definition.} \(e=(0,1,v)^{\mathsf T}\)
  \item \texttt{f} --- \texttt{vector}; \(\mathbb R^3\)
  \par\smallskip\noindent\emph{Definition.} \(f=(0,0,1)^{\mathsf T}\)
  \item \texttt{k} --- \texttt{vector}; \(\mathbb R^3\)
  \par\smallskip\noindent\emph{Definition.} \(k=g\times d=(-uz,z+vy,-y)^{\mathsf T}\)
  \item \texttt{\textbackslash{}\allowbreak{}widetilde k} --- \texttt{vector}; \(\mathbb R^3\)
  \par\smallskip\noindent\emph{Definition.} a nonzero representative of \(\ker(M_a)\)
  \item \texttt{N} --- \texttt{vector}; \(\mathbb R^3\)
  \par\smallskip\noindent\emph{Definition.} \(N=\operatorname{adj}(M_a)e_c\)
  \item \texttt{W} --- \texttt{vector}; \(\mathbb R^3\)
  \par\smallskip\noindent\emph{Definition.} \(W=M_bN\)
  \item \texttt{D} --- \texttt{scalar}
  \par\smallskip\noindent\emph{Definition.} \(D=N_bW_b+N_cW_c\)
  \item \texttt{\textbackslash{}\allowbreak{}operatorname\{Rec\}\allowbreak{}} --- \texttt{partial map}; \(\operatorname{Sym}^3(\mathbb R^3)\to\mathbb R^{3\times3}\); \texttt{piecewise-\allowbreak{}rational cubic recovery map}
  \item \texttt{\textbackslash{}\allowbreak{}mathcal P\_\allowbreak{}+\allowbreak{}} --- \texttt{class of probability laws}
  \item \texttt{\textbackslash{}\allowbreak{}theta} --- \texttt{moment parameter}
  \par\smallskip\noindent\emph{Definition.} \(\theta=(\pi,u,v,\gamma,s,\kappa)\)
  \item \texttt{\textbackslash{}\allowbreak{}Theta\_\allowbreak{}+\allowbreak{}} --- \texttt{compact parameter set}
  \item \texttt{\textbackslash{}\allowbreak{}Phi\_\allowbreak{}+\allowbreak{}} --- \texttt{polynomial map}; \(\Theta_+\to\operatorname{Sym}^3(\mathbb R^3)\)
  \par\smallskip\noindent\emph{Definition.} \(\Phi_+(\theta)=hg^{\otimes3}+rd^{\otimes3}+te^{\otimes3}+wf^{\otimes3}\)
  \item \texttt{n} --- \texttt{positive integer}; \(\mathbb N\)
  \item \texttt{\textbackslash{}\allowbreak{}alpha} --- \texttt{scalar}; \((0,1)\)
  \item \texttt{X\textasciicircum{}\{(\textbackslash{}\allowbreak{}ell)\}\allowbreak{}} --- random vector indexed by \(1\le\ell\le n\); \(\mathbb R^3\)
  \item \texttt{\textbackslash{}\allowbreak{}widehat C} --- \texttt{symmetric tensor}; \(\operatorname{Sym}^3(\mathbb R^3)\); \texttt{empirical cubic tensor}
  \item \texttt{R\_\allowbreak{}X} --- \texttt{positive scalar}
  \par\smallskip\noindent\emph{Definition.} \(R_X=37/4\)
  \item \texttt{K\_\allowbreak{}0} --- \texttt{positive scalar}
  \par\smallskip\noindent\emph{Definition.} \(K_0=R_X^3\)
  \item \texttt{\textbackslash{}\allowbreak{}varepsilon\_\allowbreak{}n} --- \texttt{positive scalar}
  \par\smallskip\noindent\emph{Definition.} \(\varepsilon_n(\alpha)=R_X^3\sqrt{2\log(20/\alpha)/n}\)
  \item \texttt{\textbackslash{}\allowbreak{}mathsf\{CS\}\allowbreak{}\_\allowbreak{}n} --- \texttt{random compact set}; subsets of \(\mathbb R^{3\times3}\)
  \item \texttt{d\_\allowbreak{}W} --- \texttt{positive scalar}
  \par\smallskip\noindent\emph{Definition.} \(d_W=729/2097152\)
  \item \texttt{d\_\allowbreak{}D} --- \texttt{positive scalar}
  \par\smallskip\noindent\emph{Definition.} \(d_D=19683/34359738368\)
  \item \texttt{d\_\allowbreak{}b} --- \texttt{positive scalar}
  \par\smallskip\noindent\emph{Definition.} \(d_b=3645/33554432\)
  \item \texttt{d\_\allowbreak{}c} --- \texttt{positive scalar}
  \par\smallskip\noindent\emph{Definition.} \(d_c=19683/134217728\)
  \item \texttt{m} --- \texttt{positive scalar}
  \par\smallskip\noindent\emph{Definition.} \(m=27/65536\)
  \item \texttt{\textbackslash{}\allowbreak{}delta\_\allowbreak{}0} --- \texttt{positive scalar}
  \par\smallskip\noindent\emph{Definition.} \(\delta_0=\min\{\min(d_b,d_c)/(18K_0^2),m/(2K_0)\}\)
  \item \texttt{L\_\allowbreak{}\textbackslash{}\allowbreak{}star} --- \texttt{positive scalar}
  \par\smallskip\noindent\emph{Definition.} \(L_\star=\max\{63K_0^2/(2d_W),150K_0^4/d_D\}\)
  \item \texttt{\textbackslash{}\allowbreak{}Sigma} --- \texttt{matrix}; \(\operatorname{Sym}_{++}^3\)
  \par\smallskip\noindent\emph{Definition.} \(\Sigma=\operatorname{Cov}_P(X)\)
  \item \texttt{\textbackslash{}\allowbreak{}Psi\_\allowbreak{}\textbackslash{}\allowbreak{}pi} --- \texttt{covariance map}
  \par\smallskip\noindent\emph{Definition.} \(\Psi_\pi(u,v,\gamma,s)=A_\pi(\gamma\gamma^{\mathsf T}+\operatorname{diag}(s^2))A_\pi^{\mathsf T}\)
  \item \texttt{\textbackslash{}\allowbreak{}mathcal S\_\allowbreak{}\textbackslash{}\allowbreak{}pi} --- subset of \(\operatorname{Sym}_{++}^3\)
  \par\smallskip\noindent\emph{Definition.} \(\mathcal S_\pi=\{\Psi_\pi(u,v,\gamma,s):(u,v,\gamma_1,\gamma_2,\gamma_3)\in[1/4,3/4]^5,\ s\in[3/4,5/4]^3\}\)
  \item \texttt{\textbackslash{}\allowbreak{}Sigma\_\allowbreak{}0} --- \texttt{matrix}; \(\operatorname{Sym}_{++}^3\)
  \item \texttt{\textbackslash{}\allowbreak{}beta} --- \texttt{scalar}; \(\mathbb R\); \texttt{the published treatment-\allowbreak{}to-\allowbreak{}outcome direct effect}
  \item \texttt{\textbackslash{}\allowbreak{}ell\_\allowbreak{}\{23\}\allowbreak{}\textasciicircum{}\{\textbackslash{}\allowbreak{}mathrm\{M2\}\allowbreak{}\}\allowbreak{}} --- \texttt{polynomial variable}
  \par\smallskip\noindent\emph{Definition.} the retained released-code variable, identified by \(\ell_{23}^{\mathrm{M2}}=\beta\)
  \item \texttt{\textbackslash{}\allowbreak{}omega\textasciicircum{}\{(2)\}\allowbreak{}} --- \texttt{vector}; \(\mathbb R^4\); \texttt{unrestricted diagonal source second cumulants}
  \item \texttt{\textbackslash{}\allowbreak{}omega\textasciicircum{}\{(3)\}\allowbreak{}} --- \texttt{vector}; \(\mathbb R^4\); \texttt{unrestricted diagonal source third cumulants}
  \item \texttt{L\_\allowbreak{}\{\textbackslash{}\allowbreak{}mathrm\{pub\}\allowbreak{}\}\allowbreak{}} --- \texttt{residual transformation matrix}; \(\mathbb R^{4\times4}\)
  \par\smallskip\noindent\emph{Definition.} \(L_{\mathrm{pub}}=\begin{pmatrix}1&-\gamma_1&-\gamma_2&-\gamma_3\\0&1&-u&0\\0&0&1&-\beta\\0&0&0&1\end{pmatrix}\), the upper-triangular residual transformation with negative edge entries in the released row-vector convention
  \item \texttt{\textbackslash{}\allowbreak{}Lambda\_\allowbreak{}\{\textbackslash{}\allowbreak{}mathrm\{pub\}\allowbreak{}\}\allowbreak{}} --- \texttt{mixing matrix}; \(\mathbb R^{3\times4}\)
  \par\smallskip\noindent\emph{Definition.} the observed-row block of \(L_{\mathrm{pub}}^{-\mathsf T}\), namely \(\Lambda_{\mathrm{pub}}=\begin{pmatrix}\gamma_1&1&0&0\\u\gamma_1+\gamma_2&u&1&0\\\beta(u\gamma_1+\gamma_2)+\gamma_3&\beta u&\beta&1\end{pmatrix}\), with rows indexed by observed nodes \(1,2,3\) and columns indexed by sources \(0,1,2,3\)
  \item \texttt{\textbackslash{}\allowbreak{}mathbf m\textasciicircum{}\{(2)\}\allowbreak{}} --- \texttt{second-\allowbreak{}moment coordinate vector}; \(\mathbb R^6\)
  \par\smallskip\noindent\emph{Definition.} \(\mathbf m^{(2)}=(m_{11},m_{12},m_{13},m_{22},m_{23},m_{33})=(\Sigma_{11},\Sigma_{12},\Sigma_{13},\Sigma_{22},\Sigma_{23},\Sigma_{33})\)
  \item \texttt{\textbackslash{}\allowbreak{}mathbf c} --- \texttt{third-\allowbreak{}moment coordinate vector}; \(\mathbb R^{10}\)
  \par\smallskip\noindent\emph{Definition.} \(\mathbf c=(c_{111},c_{112},c_{113},c_{122},c_{123},c_{133},c_{222},c_{223},c_{233},c_{333})\), where \(c_{ijk}=C_{ijk}\) for \(1\le i\le j\le k\le3\)
  \item \texttt{\textbackslash{}\allowbreak{}mathbf m} --- \texttt{moment-\allowbreak{}coordinate vector}; \(\mathbb R^{16}\)
  \par\smallskip\noindent\emph{Definition.} \(\mathbf m=(\mathbf m^{(2)},\mathbf c)\)
  \item \texttt{\textbackslash{}\allowbreak{}vartheta\_\allowbreak{}\{\textbackslash{}\allowbreak{}mathbb R\}\allowbreak{}} --- \texttt{unrestricted real parameter vector}; \(\mathbb R^{13}\); \texttt{parameters of the fixed published one-\allowbreak{}proxy graph through order three}
  \par\smallskip\noindent\emph{Definition.} \(\vartheta_{\mathbb R}=(\gamma_1,\gamma_2,\gamma_3,u,\beta,\omega^{(2)},\omega^{(3)})\)
  \item \texttt{\textbackslash{}\allowbreak{}widetilde\textbackslash{}\allowbreak{}Phi\_\allowbreak{}\{\textbackslash{}\allowbreak{}le3\}\allowbreak{}} --- \texttt{polynomial map}
  \par\smallskip\noindent\emph{Definition.} \(\widetilde\Phi_{\le3}(\vartheta_{\mathbb R})=(\beta,\mathbf m)\) with \(m_{ik}=\sum_{j=0}^3\omega_j^{(2)}(\Lambda_{\mathrm{pub}})_{ij}(\Lambda_{\mathrm{pub}})_{kj}\) for \(1\le i\le k\le3\), and \(c_{ik\ell}=\sum_{j=0}^3\omega_j^{(3)}(\Lambda_{\mathrm{pub}})_{ij}(\Lambda_{\mathrm{pub}})_{kj}(\Lambda_{\mathrm{pub}})_{\ell j}\) for \(1\le i\le k\le\ell\le3\)
  \item \texttt{\textbackslash{}\allowbreak{}mathbb F} --- \texttt{coefficient field}
  \par\smallskip\noindent\emph{Definition.} \(\mathbb F=\mathbb Q\)
  \item \texttt{\textbackslash{}\allowbreak{}mathcal R\_\allowbreak{}\{\textbackslash{}\allowbreak{}mathrm\{pub\}\allowbreak{}\}\allowbreak{}} --- \texttt{polynomial ring}
  \par\smallskip\noindent\emph{Definition.} \(\mathcal R_{\mathrm{pub}}=\mathbb F[\beta,\mathbf m,\gamma_1,\gamma_2,\gamma_3,u,\omega^{(2)}_0,\ldots,\omega^{(2)}_3,\omega^{(3)}_0,\ldots,\omega^{(3)}_3]\)
  \item \texttt{\textbackslash{}\allowbreak{}mathcal J\_\allowbreak{}\{\textbackslash{}\allowbreak{}le3\}\allowbreak{}} --- \texttt{full parameter-\allowbreak{}and-\allowbreak{}moment ideal}
  \item \texttt{\textbackslash{}\allowbreak{}mathcal I\_\allowbreak{}\{\textbackslash{}\allowbreak{}le3\}\allowbreak{}} --- \texttt{elimination ideal}
  \item \texttt{\textbackslash{}\allowbreak{}mathcal I\_\allowbreak{}\{\textbackslash{}\allowbreak{}mathrm\{obs\}\allowbreak{}\}\allowbreak{}} --- \texttt{observable-\allowbreak{}image ideal}
  \par\smallskip\noindent\emph{Definition.} \(\mathcal I_{\mathrm{obs}}=\mathcal I_{\le3}\cap\mathbb F[\mathbf m]\)
  \item \texttt{K\_\allowbreak{}\{\textbackslash{}\allowbreak{}mathrm\{obs\}\allowbreak{}\}\allowbreak{}} --- \texttt{field}
  \par\smallskip\noindent\emph{Definition.} \(K_{\mathrm{obs}}=\operatorname{Frac}(\mathbb F[\mathbf m]/\mathcal I_{\mathrm{obs}})\)
  \item \texttt{N\textasciicircum{}\{\textbackslash{}\allowbreak{}mathrm\{el\}\allowbreak{}\}\allowbreak{}} --- \texttt{vector}; \(\mathbb R^3\)
  \par\smallskip\noindent\emph{Definition.} \(N^{\mathrm{el}}=(c_{112}c_{123}-c_{113}c_{122},c_{112}c_{113}-c_{111}c_{123},c_{111}c_{122}-c_{112}^2)^{\mathsf T}\)
  \item \texttt{W\_\allowbreak{}2\textasciicircum{}\{\textbackslash{}\allowbreak{}mathrm\{el\}\allowbreak{}\}\allowbreak{}} --- \texttt{polynomial}
  \par\smallskip\noindent\emph{Definition.} \(W_2^{\mathrm{el}}=c_{122}N^{\mathrm{el}}_1+c_{222}N^{\mathrm{el}}_2+c_{223}N^{\mathrm{el}}_3\)
  \item \texttt{W\_\allowbreak{}3\textasciicircum{}\{\textbackslash{}\allowbreak{}mathrm\{el\}\allowbreak{}\}\allowbreak{}} --- \texttt{polynomial}
  \par\smallskip\noindent\emph{Definition.} \(W_3^{\mathrm{el}}=c_{123}N^{\mathrm{el}}_1+c_{223}N^{\mathrm{el}}_2+c_{233}N^{\mathrm{el}}_3\)
  \item \texttt{R\_\allowbreak{}\{\textbackslash{}\allowbreak{}mathrm\{el\}\allowbreak{}\}\allowbreak{}} --- \texttt{polynomial}
  \par\smallskip\noindent\emph{Definition.} \(R_{\mathrm{el}}(\beta,\mathbf c)=\beta W_2^{\mathrm{el}}-W_3^{\mathrm{el}}\)
  \item \texttt{S\_\allowbreak{}W} --- \texttt{multiplicative set}
  \par\smallskip\noindent\emph{Definition.} \(S_W=\{(W_2^{\mathrm{el}})^k:k\in\mathbb N_0\}\subset\mathbb F[\beta,\mathbf m]\)
  \item \texttt{\textbackslash{}\allowbreak{}mathcal I\_\allowbreak{}\{\textbackslash{}\allowbreak{}le3\}\allowbreak{}\textasciicircum{}\{(W)\}\allowbreak{}} --- \texttt{localized elimination ideal}
  \par\smallskip\noindent\emph{Definition.} \(\mathcal I_{\le3}^{(W)}=S_W^{-1}\mathcal I_{\le3}\subset S_W^{-1}\mathbb F[\beta,\mathbf m]\); equivalently its contraction is \(\mathcal I_{\le3}:(W_2^{\mathrm{el}})^\infty\)
  \item \texttt{\textbackslash{}\allowbreak{}mathfrak A} --- \texttt{finite semialgebraic fiber atlas}
  \item \texttt{\textbackslash{}\allowbreak{}Theta\_\allowbreak{}\textbackslash{}\allowbreak{}pi\textasciicircum{}\textbackslash{}\allowbreak{}circ} --- principal semialgebraic parameter chart indexed by \(\pi\in\mathfrak S_3\)
  \item \texttt{C\_\allowbreak{}\textbackslash{}\allowbreak{}pi\textasciicircum{}\{(3)\}\allowbreak{}} --- polynomial cubic-tensor map indexed by \(\pi\in\mathfrak S_3\); \(\mathbb R^9\to\operatorname{Sym}^3(\mathbb R^3)\)
  \item \texttt{\textbackslash{}\allowbreak{}mathcal C\textasciicircum{}\textbackslash{}\allowbreak{}circ} --- \texttt{principal unknown-\allowbreak{}order cubic tensor model}; \(\operatorname{Sym}^3(\mathbb R^3)\)
\end{itemize}

\section{Assumptions}

\begin{assumption}[ass:structural-equation]\label{ass:structural-equation}
\[X=A_\pi(\gamma H+\operatorname{diag}(s)\xi)\]
\par\smallskip\noindent\emph{Standard} (canonical linear non-Gaussian structural equation, \cite{tramontanoEtAl2025HigherOrderCumulants}).
\end{assumption}

\begin{assumption}[ass:source-independence]\label{ass:source-independence}
\[H\mathrel{\perp\!\!\!\perp}\xi_1\mathrel{\perp\!\!\!\perp}\xi_2\mathrel{\perp\!\!\!\perp}\xi_3\]
\par\smallskip\noindent\emph{Standard} (independent LiNGAM sources, \cite{shimizuEtAl2006LiNGAM}).
\end{assumption}

\begin{assumption}[ass:source-centering]\label{ass:source-centering}
\[\mathbb E[(H,\xi^{\mathsf T})^{\mathsf T}]=0_4\]
\par\smallskip\noindent\emph{Standard} (centered-source normalization, \cite{caiEtAl2023HigherOrderCumulants}).
\end{assumption}

\begin{assumption}[ass:source-unit-variance]\label{ass:source-unit-variance}
\[\operatorname{diag}\operatorname{Cov}((H,\xi^{\mathsf T})^{\mathsf T})=\mathbf 1_4\]
\par\smallskip\noindent\emph{Standard} (unit-variance source normalization, \cite{caiEtAl2023HigherOrderCumulants}).
\end{assumption}

\begin{assumption}[ass:source-support]\label{ass:source-support}
\[\mathbb P((H,\xi^{\mathsf T})^{\mathsf T}\in[-2,2]^4)=1\]
\par\smallskip\noindent\emph{Standard} (bounded-support concentration condition, \cite{hoeffding1963ProbabilityInequalities}).
\end{assumption}

\begin{assumption}[ass:positive-mechanism-box]\label{ass:positive-mechanism-box}
\[(u,v,\gamma_1,\gamma_2,\gamma_3)\in[1/4,3/4]^5\]
\par\smallskip\noindent\emph{Novel.} This sign-and-margin chamber represents reinforcing propagation and a common positive latent shock, as in monotone production, demand, and biological signaling chains; it is stated as the substantive semialgebraic restriction whose identifying content is characterized.
\end{assumption}

\begin{assumption}[ass:residual-scale-box]\label{ass:residual-scale-box}
\[s\in[3/4,5/4]^3\]
\par\smallskip\noindent\emph{Novel.} The normalization permits heterogeneous idiosyncratic scales around the latent-source variance benchmark and describes a compact calibration class rather than an oracle restriction.
\end{assumption}

\begin{assumption}[ass:positive-skew-box]\label{ass:positive-skew-box}
\[\kappa\in[1/2,1]^4\]
\par\smallskip\noindent\emph{Novel.} A common skew direction arises for centered bounded arrival, demand, and shock variables; the interval permits unequal source laws and rules out only weak or sign-canceling skewness.
\end{assumption}

\begin{assumption}[ass:iid-sampling]\label{ass:iid-sampling}
\[(X^{(1)},\ldots,X^{(n)})\sim P^{\otimes n}\]
\par\smallskip\noindent\emph{Standard} (independent and identically distributed sampling, \cite{tramontanoEtAl2025HigherOrderCumulants}).
\end{assumption}

\section{Definitions}

\begin{definition}[def:positive-model]\label{def:positive-model}
\par\smallskip\noindent\emph{Defined object.} \texttt{\textbackslash{}\allowbreak{}mathcal P\_\allowbreak{}+\allowbreak{}}
\par\smallskip\noindent\emph{Construction.} \[\mathcal P_+=\{P:\exists(\pi,u,v,\gamma,s,H,\xi)\text{ satisfying }X=A_\pi(\gamma H+\operatorname{diag}(s)\xi),\ H\perp\!\!\!\perp\xi_1\perp\!\!\!\perp\xi_2\perp\!\!\!\perp\xi_3,\ \mathbb E[(H,\xi^{\mathsf T})^{\mathsf T}]=0_4,\ \operatorname{diag}\operatorname{Cov}((H,\xi^{\mathsf T})^{\mathsf T})=\mathbf1_4,\ \mathbb P((H,\xi^{\mathsf T})^{\mathsf T}\in[-2,2]^4)=1,\ (u,v,\gamma_1,\gamma_2,\gamma_3)\in[1/4,3/4]^5,\ s\in[3/4,5/4]^3,\ \kappa\in[1/2,1]^4\}\]
\par\smallskip\noindent\emph{Member properties.} \texttt{ass:\allowbreak{}structural-\allowbreak{}equation}; \texttt{ass:\allowbreak{}source-\allowbreak{}independence}; \texttt{ass:\allowbreak{}source-\allowbreak{}centering}; \texttt{ass:\allowbreak{}source-\allowbreak{}unit-\allowbreak{}variance}; \texttt{ass:\allowbreak{}source-\allowbreak{}support}; \texttt{ass:\allowbreak{}positive-\allowbreak{}mechanism-\allowbreak{}box}; \texttt{ass:\allowbreak{}residual-\allowbreak{}scale-\allowbreak{}box}; \texttt{ass:\allowbreak{}positive-\allowbreak{}skew-\allowbreak{}box}.
\end{definition}

\begin{definition}[def:positive-parameter-set]\label{def:positive-parameter-set}
\par\smallskip\noindent\emph{Defined object.} \texttt{\textbackslash{}\allowbreak{}Theta\_\allowbreak{}+\allowbreak{}}
\par\smallskip\noindent\emph{Construction.} \[\Theta_+=\mathfrak S_3\times[1/4,3/4]^2\times[1/4,3/4]^3\times[3/4,5/4]^3\times[1/2,1]^4\]
\par\smallskip\noindent\emph{Inputs.} \texttt{ass:\allowbreak{}positive-\allowbreak{}mechanism-\allowbreak{}box}; \texttt{ass:\allowbreak{}residual-\allowbreak{}scale-\allowbreak{}box}; \texttt{ass:\allowbreak{}positive-\allowbreak{}skew-\allowbreak{}box}.
\end{definition}

\begin{definition}[def:cubic-recovery]\label{def:cubic-recovery}
\par\smallskip\noindent\emph{Defined object.} \texttt{\textbackslash{}\allowbreak{}operatorname\{Rec\}\allowbreak{}}
\par\smallskip\noindent\emph{Construction.} On the certified principal domain \(\mathcal C^{\circ}\), define \(\operatorname{Rec}\) by the following partial rational rule.  For \(C\in\mathcal C^{\circ}\), set
\[
a=\operatorname*{arg\,unique}_{i\in\mathcal V}\{\operatorname{rank}(M_i)=2\},
\qquad
\widetilde k\in\ker(M_a)\setminus\{0\},
\]
\[
b=\operatorname*{arg\,unique}_{j\ne a}
\{\operatorname{sgn}(\widetilde k_j)=-\operatorname{sgn}(\widetilde k_a)\},
\qquad
\{c\}=\mathcal V\setminus\{a,b\},
\qquad
\pi=(a,b,c),
\]
and then set
\[
N=\operatorname{adj}(M_a)e_c,
\qquad
W=M_bN,
\qquad
D=N_bW_b+N_cW_c,
\qquad
\operatorname{Rec}(C)
=B_{\pi}\!\left(-\frac{N_aW_b}{D},\frac{W_c}{W_b}\right).
\]
The principal-chart theorem proves that every uniqueness assertion in this rule holds and that \(W_bD\ne0\) on \(\mathcal C^{\circ}\).  If \(C\notin\mathcal C^{\circ}\), this displayed principal branch is left undefined: neither an alternate rational chart nor fiber multiplicity is decided by the definition.
\par\smallskip\noindent\emph{Inputs.} \texttt{def:\allowbreak{}principal-\allowbreak{}real-\allowbreak{}cubic-\allowbreak{}chart}.
\end{definition}

\begin{definition}[def:empirical-cubic]\label{def:empirical-cubic}
\par\smallskip\noindent\emph{Defined object.} \texttt{\textbackslash{}\allowbreak{}widehat C}
\par\smallskip\noindent\emph{Construction.} \[\widehat C_{ijk}=n^{-1}\sum_{\ell=1}^nX_i^{(\ell)}X_j^{(\ell)}X_k^{(\ell)},\qquad 1\le i\le j\le k\le3\]
\par\smallskip\noindent\emph{Inputs.} \texttt{ass:\allowbreak{}source-\allowbreak{}centering}; \texttt{ass:\allowbreak{}iid-\allowbreak{}sampling}.
\end{definition}

\begin{definition}[def:confidence-set]\label{def:confidence-set}
\par\smallskip\noindent\emph{Defined object.} \texttt{\textbackslash{}\allowbreak{}mathsf\{CS\}\allowbreak{}\_\allowbreak{}n}
\par\smallskip\noindent\emph{Construction.} \[\mathsf{CS}_n(\alpha)=\{B_\pi(u,v):\theta=(\pi,u,v,\gamma,s,\kappa)\in\Theta_+,\ \|\Phi_+(\theta)-\widehat C\|_\infty\le\varepsilon_n(\alpha)\}\]
\par\smallskip\noindent\emph{Inputs.} \texttt{def:\allowbreak{}positive-\allowbreak{}parameter-\allowbreak{}set}; \texttt{def:\allowbreak{}empirical-\allowbreak{}cubic}; \texttt{ass:\allowbreak{}iid-\allowbreak{}sampling}.
\end{definition}

\begin{definition}[def:covariance-witness]\label{def:covariance-witness}
\par\smallskip\noindent\emph{Defined object.} \texttt{\textbackslash{}\allowbreak{}Sigma\_\allowbreak{}0}
\par\smallskip\noindent\emph{Construction.} \[\Sigma_0=\begin{pmatrix}179/100&817/1000&4507/6100\\817/1000&12891/10000&817/1000\\4507/6100&817/1000&179/100\end{pmatrix}\]
\par\smallskip\noindent\emph{Inputs.} \texttt{ass:\allowbreak{}source-\allowbreak{}unit-\allowbreak{}variance}; \texttt{ass:\allowbreak{}positive-\allowbreak{}mechanism-\allowbreak{}box}; \texttt{ass:\allowbreak{}residual-\allowbreak{}scale-\allowbreak{}box}.
\end{definition}

\begin{definition}[def:published-mixing-map]\label{def:published-mixing-map}
\par\smallskip\noindent\emph{Defined object.} \texttt{\textbackslash{}\allowbreak{}Lambda\_\allowbreak{}\{\textbackslash{}\allowbreak{}mathrm\{pub\}\allowbreak{}\}\allowbreak{}}
\par\smallskip\noindent\emph{Construction.} For the exact graph with latent node \(0\), observed nodes \(1,2,3\), and edges \(0\to1\), \(0\to2\), \(0\to3\), \(1\to2\), and \(2\to3\), use the released row-vector residual transform \[L_{\mathrm{pub}}=\begin{pmatrix}1&-\gamma_1&-\gamma_2&-\gamma_3\\0&1&-u&0\\0&0&1&-\beta\\0&0&0&1\end{pmatrix},\qquad \Lambda_{\mathrm{pub}}=\begin{pmatrix}\gamma_1&1&0&0\\u\gamma_1+\gamma_2&u&1&0\\\beta(u\gamma_1+\gamma_2)+\gamma_3&\beta u&\beta&1\end{pmatrix}.\] The latter is the observed-row block of \(L_{\mathrm{pub}}^{-\mathsf T}\). Define \(m_{ik}=\sum_{j=0}^3\omega_j^{(2)}(\Lambda_{\mathrm{pub}})_{ij}(\Lambda_{\mathrm{pub}})_{kj}\) and \(c_{ik\ell}=\sum_{j=0}^3\omega_j^{(3)}(\Lambda_{\mathrm{pub}})_{ij}(\Lambda_{\mathrm{pub}})_{kj}(\Lambda_{\mathrm{pub}})_{\ell j}\). The retained coordinates are exactly \(\mathbf m^{(2)}=(m_{11},m_{12},m_{13},m_{22},m_{23},m_{33})\) and \(\mathbf c=(c_{111},c_{112},c_{113},c_{122},c_{123},c_{133},c_{222},c_{223},c_{233},c_{333})\).
\end{definition}

\begin{definition}[def:full-published-moment-ideal]\label{def:full-published-moment-ideal}
\par\smallskip\noindent\emph{Defined object.} \texttt{\textbackslash{}\allowbreak{}mathcal J\_\allowbreak{}\{\textbackslash{}\allowbreak{}le3\}\allowbreak{}}
\par\smallskip\noindent\emph{Construction.} \[\mathcal J_{\le3}=\left\langle m_{ik}-\sum_{j=0}^3\omega_j^{(2)}(\Lambda_{\mathrm{pub}})_{ij}(\Lambda_{\mathrm{pub}})_{kj}:1\le i\le k\le3\right\rangle+\left\langle c_{ik\ell}-\sum_{j=0}^3\omega_j^{(3)}(\Lambda_{\mathrm{pub}})_{ij}(\Lambda_{\mathrm{pub}})_{kj}(\Lambda_{\mathrm{pub}})_{\ell j}:1\le i\le k\le\ell\le3\right\rangle\subset\mathcal R_{\mathrm{pub}}.\] This is the parameterized form of the released residual-diagonal covariance equations and off-diagonal cubic residual equations over \(\mathbb F=\mathbb Q\).
\par\smallskip\noindent\emph{Inputs.} \texttt{def:\allowbreak{}published-\allowbreak{}mixing-\allowbreak{}map}.
\end{definition}

\begin{definition}[def:published-elimination-ideal]\label{def:published-elimination-ideal}
\par\smallskip\noindent\emph{Defined object.} \texttt{\textbackslash{}\allowbreak{}mathcal I\_\allowbreak{}\{\textbackslash{}\allowbreak{}le3\}\allowbreak{}}
\par\smallskip\noindent\emph{Construction.} \[\mathcal I_{\le3}=\mathcal J_{\le3}\cap\mathbb F[\beta,\mathbf m],\qquad \mathcal I_{\mathrm{obs}}=\mathcal I_{\le3}\cap\mathbb F[\mathbf m],\qquad K_{\mathrm{obs}}=\operatorname{Frac}(\mathbb F[\mathbf m]/\mathcal I_{\mathrm{obs}}).\] Thus \(\mathcal I_{\le3}\) eliminates \(\gamma_1,\gamma_2,\gamma_3,u,\omega^{(2)},\omega^{(3)}\) while retaining the published effect \(\beta\) and all sixteen observed moment coordinates.
\par\smallskip\noindent\emph{Inputs.} \texttt{def:\allowbreak{}full-\allowbreak{}published-\allowbreak{}moment-\allowbreak{}ideal}.
\end{definition}

\begin{definition}[def:published-localization]\label{def:published-localization}
\par\smallskip\noindent\emph{Defined object.} \texttt{\textbackslash{}\allowbreak{}mathcal I\_\allowbreak{}\{\textbackslash{}\allowbreak{}le3\}\allowbreak{}\textasciicircum{}\{(W)\}\allowbreak{}}
\par\smallskip\noindent\emph{Construction.} \[S_W=\{(W_2^{\mathrm{el}})^k:k\in\mathbb N_0\},\qquad \mathcal I_{\le3}^{(W)}=S_W^{-1}\mathcal I_{\le3}\subset S_W^{-1}\mathbb F[\beta,\mathbf m].\] Localization therefore means inversion of the single observable polynomial \(W_2^{\mathrm{el}}\); its contraction to \(\mathbb F[\beta,\mathbf m]\) is the saturation \(\mathcal I_{\le3}:(W_2^{\mathrm{el}})^\infty\). Extension afterward to \(K_{\mathrm{obs}}[\beta]\) inverts every nonzero observable class.
\par\smallskip\noindent\emph{Inputs.} \texttt{def:\allowbreak{}published-\allowbreak{}elimination-\allowbreak{}ideal}.
\end{definition}

\begin{definition}[def:linear-elimination-certificate]\label{def:linear-elimination-certificate}
\par\smallskip\noindent\emph{Defined object.} \texttt{R\_\allowbreak{}\{\textbackslash{}\allowbreak{}mathrm\{el\}\allowbreak{}\}\allowbreak{}}
\par\smallskip\noindent\emph{Construction.} \[R_{\mathrm{el}}(\beta,\mathbf c)=\beta(c_{122}N^{\mathrm{el}}_1+c_{222}N^{\mathrm{el}}_2+c_{223}N^{\mathrm{el}}_3)-(c_{123}N^{\mathrm{el}}_1+c_{223}N^{\mathrm{el}}_2+c_{233}N^{\mathrm{el}}_3)\]
\par\smallskip\noindent\emph{Inputs.} \texttt{def:\allowbreak{}published-\allowbreak{}elimination-\allowbreak{}ideal}.
\end{definition}

\begin{definition}[def:atlas-handle]\label{def:atlas-handle}
\par\smallskip\noindent\emph{Defined object.} \texttt{\textbackslash{}\allowbreak{}mathfrak A}
\par\smallskip\noindent\emph{Construction.} Apply chartwise saturation of the full covariance-and-cubic moment ideal by every nonzero \(2\times2\) slice minor and every adjugate contraction denominator, then perform real primary decomposition and solve the six order-specific fibers on each sign cell.
\par\smallskip\noindent\emph{Inputs.} \texttt{def:\allowbreak{}published-\allowbreak{}elimination-\allowbreak{}ideal}; \texttt{def:\allowbreak{}cubic-\allowbreak{}recovery}.
\end{definition}

\begin{definition}[def:principal-real-cubic-chart]\label{def:principal-real-cubic-chart}
\par\smallskip\noindent\emph{Defined object.} \texttt{\textbackslash{}\allowbreak{}mathcal C\textasciicircum{}\textbackslash{}\allowbreak{}circ}
\par\smallskip\noindent\emph{Construction.} For each fixed labeled order \(\pi=(a,b,c)\in\mathfrak S_3\) and each \((h,r,t,w,x,y,z,u,v)\in\mathbb R^9\), put \(p=ux+y\) and \(q=vp+z\).  Define the cubic polynomial map, in the ambient labeled coordinates, by
\[
C_{\pi}^{(3)}(h,r,t,w,x,y,z,u,v)
=h(xe_a+pe_b+qe_c)^{\otimes3}
+r(e_a+ue_b+uve_c)^{\otimes3}
+t(e_b+ve_c)^{\otimes3}
+we_c^{\otimes3}.
\]
Its principal real parameter chart is the semialgebraic set
\[
\Theta_{\pi}^{\circ}
=\left\{(h,r,t,w,x,y,z,u,v)\in\mathbb R^9:
h r t x y z u p\,\det\!\left(C_{\pi}^{(3)}(h,r,t,w,x,y,z,u,v)(e_c,\cdot,\cdot)\right)\ne0,
\quad uz(z+vy)>0,
\quad uzy>0
\right\}.
\]
Equivalently, the first condition is \(h r t x y z u p\det(M_c)\ne0\) for the tensor displayed above.  Define the unknown-order principal tensor model by
\[
\mathcal C^{\circ}
=\bigcup_{\pi\in\mathfrak S_3}C_{\pi}^{(3)}(\Theta_{\pi}^{\circ})
\subset\operatorname{Sym}^3(\mathbb R^3).
\]
No positivity, boundedness, probability-law realizability, or restriction on \(w\) or \(v\) is imposed beyond the displayed conditions.
\end{definition}

\section{Main results}

\begin{proposition}[prop:source-feasibility]\label{prop:source-feasibility}
For every \(\kappa_0\in[1/2,1]\), let \(b_0=(\kappa_0+\sqrt{\kappa_0^2+4})/2\) and \(a_0=1/b_0\). The two-point law placing probabilities \(b_0/(a_0+b_0)\) and \(a_0/(a_0+b_0)\) on \(-a_0\) and \(b_0\), respectively, is centered, has variance one, support in \([-2,2]\), and third cumulant \(\kappa_0\). Independent products of four such laws make \(\mathcal P_+\) nonempty for every \(\kappa\in[1/2,1]^4\).
\end{proposition}
\begin{proof}
Fix \(\kappa_0\in[1/2,1]\).  The definition of \(b_0\) makes it the positive root of
\[
b_0^2-\kappa_0b_0-1=0.
\]
Consequently, using \(a_0=b_0^{-1}\),
\[
a_0b_0=1,
\qquad
b_0-a_0=\kappa_0.
\]
Let \(S\) take the values \(-a_0\) and \(b_0\) with probabilities
\[
p_-:=\frac{b_0}{a_0+b_0},
\qquad
p_+:=\frac{a_0}{a_0+b_0},
\]
respectively.  Direct calculation gives
\[
\mathbb E[S]
=\frac{-a_0b_0+b_0a_0}{a_0+b_0}=0,
\]
\[
\mathbb E[S^2]
=\frac{a_0^2b_0+b_0^2a_0}{a_0+b_0}
=a_0b_0=1,
\]
and
\[
\mathbb E[S^3]
=\frac{-a_0^3b_0+b_0^3a_0}{a_0+b_0}
=a_0b_0(b_0-a_0)
=\kappa_0.
\]
Because \(S\) is centered, its third cumulant equals \(\mathbb E[S^3]\), hence equals \(\kappa_0\).  Moreover, \(\kappa_0>0\) and the quadratic equation above imply \(b_0>1\), while
\[
b_0\le \frac{1+\sqrt5}{2}<2,
\qquad
0<a_0<1.
\]
Thus both support points belong to \([-2,2]\).

Now fix any \(\kappa=(\kappa_H,\kappa_1,\kappa_2,\kappa_3)\in[1/2,1]^4\).  Apply the preceding construction separately to these four coordinates and take the product law.  By construction the resulting variables \(H,\xi_1,\xi_2,\xi_3\) are mutually independent, centered, variance one, supported on \([-2,2]\), and have the prescribed third cumulants.  Finally, choose any \(\pi\in\mathfrak S_3\), set, for example, \(u=v=1/2\), \(\gamma=(1/2,1/2,1/2)^{\mathsf T}\), and \(s=(1,1,1)^{\mathsf T}\), and define \(X\) by the structural equation.  Every box restriction in the definition of \(\mathcal P_+\) is then satisfied.  Hence \(\mathcal P_+\) is nonempty for every prescribed \(\kappa\in[1/2,1]^4\).
\end{proof}
\par\smallskip\noindent \textit{Justification.} This prevents the compact skewness and support restrictions from describing an empty or moment-incompatible model. \textit{Gap.} The cited LiNGAM papers normalize source moments but do not verify simultaneous bounded-support realizability for the entire maintained skewness interval. \textit{Consumer.} The finite-sample inversion theorem uses this result to interpret every retained moment parameter as an attainable source law.

\begin{lemma}[lem:cubic-factorization]\label{lem:cubic-factorization}
For every \(P\in\mathcal P_+\) with \(\pi=(a,b,c)\), the tensor in causal coordinates satisfies \[C=hg^{\otimes3}+rd^{\otimes3}+te^{\otimes3}+wf^{\otimes3}.\] Its slices obey \[M_a=hxgg^{\mathsf T}+rdd^{\mathsf T},\qquad \ker(M_a)=\operatorname{span}\{k\},\] and \[\det(M_b)=hrtu p z^2,\qquad \det(M_c)=hrtu v^2qz^2+hrwuvqy^2+htwvqx^2+rtwuv^2.\]
\end{lemma}
\begin{proof}
Work in the causal coordinate order \((a,b,c)\).  Since \(B_\pi^3=0\), solving the structural equations successively gives
\[
X_a=xH+s_a\xi_a,
\]
\[
X_b=uX_a+yH+s_b\xi_b
=pH+us_a\xi_a+s_b\xi_b,
\]
and
\[
X_c=vX_b+zH+s_c\xi_c
=qH+uvs_a\xi_a+vs_b\xi_b+s_c\xi_c.
\]
Equivalently,
\[
X=gH+s_ad\xi_a+s_be\xi_b+s_cf\xi_c.
\]
The source-centering conditions contained in the definition of \(\mathcal P_+\) imply \(\mathbb E[X]=0\), so the third cumulant tensor is the third central-moment tensor \(C=\mathbb E[X^{\otimes3}]\).  Expand the last display trilinearly.  Any mixed term contains a centered source independent of the product of the remaining source factors, and therefore has expectation zero by source independence.  Only the four pure-source terms remain.  With the definitions of \(h,r,t,w\), this proves
\[
C=hg^{\otimes3}+rd^{\otimes3}+te^{\otimes3}+wf^{\otimes3}.
\]

Contracting the first tensor coordinate with the coordinate vectors gives
\[
M_a=h g_a gg^{\mathsf T}+r d_a dd^{\mathsf T}
=hxgg^{\mathsf T}+rdd^{\mathsf T},
\]
because \((g_a,d_a,e_a,f_a)=(x,1,0,0)\).  Direct expansion gives
\[
g\times d
=\begin{pmatrix}puv-qu\\q-xuv\\xu-p\end{pmatrix}
=\begin{pmatrix}-uz\\z+vy\\-y\end{pmatrix}
=k.
\]
On \(\mathcal P_+\), \(y,z>0\), so \(k\ne0\) and \(g,d\) are linearly independent.  Also \(hx>0\) and \(r>0\).  Therefore, for any \(z_0\in\mathbb R^3\),
\[
z_0^{\mathsf T}M_az_0
=hx(g^{\mathsf T}z_0)^2+r(d^{\mathsf T}z_0)^2,
\]
which vanishes exactly when \(z_0\) is orthogonal to both \(g\) and \(d\).  Hence
\[
\operatorname{rank}(M_a)=2,
\qquad
\ker(M_a)=\operatorname{span}\{g\times d\}
=\operatorname{span}\{k\}.
\]

For the middle slice,
\[
M_b=hpgg^{\mathsf T}+rudd^{\mathsf T}+tee^{\mathsf T}.
\]
Writing \(U_b=[g\ d\ e]\), determinant multiplicativity yields
\[
\det(M_b)
=\det(U_b)^2(hp)(ru)t.
\]
The defining equations \(p=ux+y\) and \(q=vp+z\) give
\[
\det[g\ d\ e]=q-vp=z,
\]
and consequently
\[
\det(M_b)=hrtu p z^2.
\]

Finally,
\[
M_c=hqgg^{\mathsf T}+ruvdd^{\mathsf T}+tvee^{\mathsf T}+wff^{\mathsf T}.
\]
Expanding its determinant trilinearly in the four rank-one summands, a term is zero whenever a direction is repeated; the four surviving terms select three distinct directions.  The needed determinants are
\[
\det[g\ d\ e]=z,
\qquad
\det[g\ d\ f]=-y,
\qquad
\det[g\ e\ f]=x,
\qquad
\det[d\ e\ f]=1.
\]
Multiplying the corresponding three weights and squaring these determinants gives
\[
\det(M_c)
=hrtu v^2qz^2+hrwuvqy^2+htwvqx^2+rtwuv^2,
\]
as claimed.
\end{proof}
\par\smallskip\noindent \textit{Justification.} The factorization exposes both the unique low-rank root slice and the contractions that isolate the two directed coefficients. \textit{Gap.} Cai et al. (2023) factor a one-latent block and Schkoda, Robeva, and Drton (2024) stack several cumulant orders; neither records this joint cubic chain factorization and its three slice determinants. \textit{Consumer.} The rank, sign, denominator, and perturbation arguments all reduce to these identities.

\begin{theorem}[thm:global-cubic-inverse]\label{thm:global-cubic-inverse}
For every \(P\in\mathcal P_+\), exactly one slice \(M_i\) has rank two, and its label is \(a\); the other two slices have rank three. Every nonzero \(\widetilde k\in\ker(M_a)\) has nonzero coordinates, and \(b\) is the unique label whose coordinate has sign opposite to \(\widetilde k_a\). With \(c=\mathcal V\setminus\{a,b\}\), \(N=\operatorname{adj}(M_a)e_c\), and \(W=M_bN\),
\[
v=\frac{W_c}{W_b},\qquad u=-\frac{N_aW_b}{N_bW_b+N_cW_c},\qquad \operatorname{Rec}(C)=B_\pi(u,v).
\]
Uniformly on \(\mathcal P_+\),
\[
|W_b|\ge d_W,\qquad |D|\ge d_D,\qquad |\det(M_b)|\ge d_b,\qquad |\det(M_c)|\ge d_c.
\]
Thus \(C_3(P)\) globally identifies the labeled chain and both direct coefficients, and \(\operatorname{Rec}\) has constant arithmetic cost.
\end{theorem}
\begin{proof}
Fix \(P\in\mathcal P_+\) and write its causal order as \(\pi=(a,b,c)\).  By the definition of \(\mathcal P_+\), the structural equation and source independence hold, and the cubic-factorization lemma therefore gives
\[
C=C_{\pi}^{(3)}(h,r,t,w,x,y,z,u,v).
\]
The positive mechanism, scale, and skewness assumptions imply
\[
h\ge\frac12,
\qquad
r,t,w\ge\left(\frac34\right)^3\frac12=\frac{27}{128},
\qquad
x,y,z,u,v\ge\frac14.
\]
In particular,
\[
p=ux+y\ge\frac{5}{16},
\qquad
q=vp+z\ge\frac{21}{64}.
\]
The principal algebra lemma gives
\[
\det(M_c)=hrtuv^2qz^2+hrwuvqy^2+htwvqx^2+rtwuv^2.
\]
Every factor in every summand is positive, so \(\det(M_c)>0\).  Moreover,
\[
hrtxyzup\det(M_c)>0,
\qquad
uz(z+vy)>0,
\qquad
uzy>0.
\]
Thus the parameter tuple lies in \(\Theta_{\pi}^{\circ}\).  The principal real cubic inverse theorem now proves the unique rank pattern, the kernel sign rule, the legitimacy of both ratios, the identities recovering \(u,v\), global uniqueness across the six labels, and
\[
\operatorname{Rec}(C)=B_{\pi}(u,v).
\]

It remains to prove the quantitative margins specific to the compact positive class.  Put \(K=hxr\).  The preceding bounds give
\[
K\ge\frac12\frac14\frac{27}{128}=\frac{27}{1024}.
\]
Using the exact contraction identities from the principal algebra lemma,
\[
|W_b|=Ktyz
\ge \frac{27}{1024}\frac{27}{128}\frac14\frac14
=\frac{729}{2097152}=d_W,
\]
and
\[
|D|=K^2ty^2z^2
\ge\left(\frac{27}{1024}\right)^2\frac{27}{128}
\left(\frac14\right)^4
=\frac{19683}{34359738368}=d_D.
\]
The determinant identity for the middle slice gives
\[
|\det(M_b)|=hrtupz^2
\ge\frac12\left(\frac{27}{128}\right)^2
\frac14\frac{5}{16}\left(\frac14\right)^2
=\frac{3645}{33554432}=d_b.
\]
Finally, every summand in \(\det(M_c)\) is positive, so retaining only the last one yields
\[
|\det(M_c)|=\det(M_c)\ge rtwuv^2
\ge\left(\frac{27}{128}\right)^3\frac14
\left(\frac14\right)^2
=\frac{19683}{134217728}=d_c.
\]
The target \(B_{\pi}(u,v)\) is structural, while the recovery map is an observable cubic functional; their equality and the denominator bounds have been derived from the stated model restrictions.  Constant arithmetic cost follows from the fixed three-variable calculation established on the principal chart.
\end{proof}
\par\smallskip\noindent \textit{Justification.} This is the theorem-level kernel: skewness alone recovers the unknown mechanism globally rather than selecting among algebraic candidates. \textit{Gap.} Tramontano et al. (2025), Theorem 3.6, reports nonidentification through order three on the fixed unrestricted graph; Cai et al. (2023) does not recover this unknown six-chain structure from \(C_3\) alone. \textit{Consumer.} CEId-from-Moments can replace its fourth-cumulant one-proxy initializer by a certified cubic branch on the positive-skew chamber.

\begin{theorem}[thm:localized-full-ideal-reaudit]\label{thm:localized-full-ideal-reaudit}
For the exact published one-proxy moment map \(\Lambda_{\mathrm{pub}}\), the observable class of \(W_2^{\mathrm{el}}\) is nonzero in \(\mathbb F[\mathbf m]/\mathcal I_{\mathrm{obs}}\), and \(R_{\mathrm{el}}=\beta W_2^{\mathrm{el}}-W_3^{\mathrm{el}}\) belongs to \(\mathcal I_{\leq3}\).  Consequently,
\[
\beta-\frac{W_3^{\mathrm{el}}}{W_2^{\mathrm{el}}}\in\mathcal I_{\leq3}^{(W)}
\qquad\text{and}\qquad
\mathcal I_{\leq3}K_{\mathrm{obs}}[\beta]
=\left\langle\beta-\frac{W_3^{\mathrm{el}}}{W_2^{\mathrm{el}}}\right\rangle.
\]
Thus every parameterized image point with \(W_2^{\mathrm{el}}\neq0\) has \(\beta=W_3^{\mathrm{el}}/W_2^{\mathrm{el}}\), and the generic full covariance-and-cubic fiber of \(\beta\) has degree one on this nonempty principal chart.  Extra roots of an individual retained eliminant are not additional common roots of the full localized ideal.  The conclusion does not classify the exceptional divisor \(W_2^{\mathrm{el}}=0\) and does not assert identification on the entire unrestricted real parameter class.
\end{theorem}
\begin{proof}
Let
\[
T=\mathbb F[\beta,\gamma_1,\gamma_2,\gamma_3,u,
\omega^{(2)}_0,\ldots,\omega^{(2)}_3,
\omega^{(3)}_0,\ldots,\omega^{(3)}_3].
\]
For each of the sixteen retained moment coordinates \(m_\nu\) in \(\mathbf m\), let \(f_\nu\in T\) denote the corresponding polynomial obtained from \(\Lambda_{\mathrm{pub}}\) in \texttt{def:published-mixing-map}.  Define the \(\mathbb F\)-algebra homomorphism
\[
\rho:\mathcal R_{\mathrm{pub}}\longrightarrow T,
\qquad
\rho(m_\nu)=f_\nu,
\]
and let \(\rho\) fix every parameter generator of \(T\).  By \texttt{def:full-published-moment-ideal},
\[
\mathcal J_{\leq3}=\langle m_\nu-f_\nu:\nu\rangle,
\]
so \(\mathcal J_{\leq3}\subseteq\ker\rho\).  Conversely, in \(\mathcal R_{\mathrm{pub}}/\mathcal J_{\leq3}\) every class \([m_\nu]\) equals \([f_\nu]\).  Hence the map induced by \(\rho\) and the map
\[
T\longrightarrow\mathcal R_{\mathrm{pub}}/\mathcal J_{\leq3}
\]
sending each parameter generator to its residue class are inverse on all generators.  Therefore
\[
\mathcal R_{\mathrm{pub}}/\mathcal J_{\leq3}\cong T,
\qquad
\ker\rho=\mathcal J_{\leq3}.
\]
Because \(T\) is a polynomial ring over the field \(\mathbb F=\mathbb Q\), it is an integral domain.  Thus \(\mathcal J_{\leq3}\) is prime.  A direct contraction argument now shows that both ideals defined in \texttt{def:published-elimination-ideal} are prime: if a product belongs to a contraction of \(\mathcal J_{\leq3}\), it belongs to \(\mathcal J_{\leq3}\), so one factor belongs to \(\mathcal J_{\leq3}\), and that factor remains in the subring being contracted.  Consequently,
\[
\mathcal I_{\leq3}=\mathcal J_{\leq3}\cap\mathbb F[\beta,\mathbf m]
\quad\text{and}\quad
\mathcal I_{\mathrm{obs}}=\mathcal I_{\leq3}\cap\mathbb F[\mathbf m]
\]
are prime, and
\[
A:=\mathbb F[\mathbf m]/\mathcal I_{\mathrm{obs}}
\]
is an integral domain with fraction field \(K_{\mathrm{obs}}\).

We next prove that the localization denominator is a nonzero observable class.  Take the rational parameter point
\[
\gamma_1=\gamma_2=\gamma_3=u=\beta=1,
\qquad
\omega_j^{(2)}=\omega_j^{(3)}=1
\quad(0\leq j\leq3),
\]
and assign the moment coordinates their values under the polynomial moment map.  This is an algebraic image point; no positivity or probabilistic realizability assertion is needed.  At this point \texttt{def:published-mixing-map} gives
\[
\Lambda_{\mathrm{pub}}=
\begin{pmatrix}
1&1&0&0\\
2&1&1&0\\
3&1&1&1
\end{pmatrix}.
\]
Its cubic moment coordinates, in the order fixed by \(\mathbf c\), are
\[
\mathbf c=(2,3,4,5,7,10,10,14,20,30).
\]
Indeed, for example, \(c_{123}=1\cdot2\cdot3+1\cdot1\cdot1=7\), and all the other displayed entries follow from the same source-column sum in \texttt{def:published-mixing-map}.  Substitution into \(N^{\mathrm{el}}\) gives
\[
N^{\mathrm{el}}
=\begin{pmatrix}
3\cdot7-4\cdot5\\
3\cdot4-2\cdot7\\
2\cdot5-3^2
\end{pmatrix}
=\begin{pmatrix}1\\-2\\1\end{pmatrix}.
\]
Therefore
\[
W_2^{\mathrm{el}}=5-20+14=-1,
\qquad
W_3^{\mathrm{el}}=7-28+20=-1.
\]
Every member of \(\mathcal I_{\mathrm{obs}}\) lies in \(\mathcal J_{\leq3}=\ker\rho\) and hence vanishes at every parameterized image point.  The nonzero evaluation above proves
\[
W_2^{\mathrm{el}}\notin\mathcal I_{\mathrm{obs}}.
\]
Thus the class of \(W_2^{\mathrm{el}}\) in \(A\) is nonzero, and its principal open chart is nonempty.

By \texttt{prop:analytic-full-ideal-membership},
\[
R_{\mathrm{el}}=\beta W_2^{\mathrm{el}}-W_3^{\mathrm{el}}
\in\mathcal I_{\leq3}.
\]
The multiplicative set in \texttt{def:published-localization} makes \(W_2^{\mathrm{el}}\) invertible.  Multiplying this membership by its inverse yields
\[
\beta-\frac{W_3^{\mathrm{el}}}{W_2^{\mathrm{el}}}
=\frac{R_{\mathrm{el}}}{W_2^{\mathrm{el}}}
\in\mathcal I_{\leq3}^{(W)}.
\]

It remains to determine the whole ideal after passage to the observable function field.  Let \(\overline{\mathcal I}\) be the image of \(\mathcal I_{\leq3}\) under the quotient map
\[
\mathbb F[\beta,\mathbf m]\longrightarrow A[\beta].
\]
Since \(\mathcal I_{\leq3}\) is prime and contains \(\mathcal I_{\mathrm{obs}}\), the quotient
\[
A[\beta]/\overline{\mathcal I}
\cong
\mathbb F[\beta,\mathbf m]/\mathcal I_{\leq3}
\]
is an integral domain.  Hence \(\overline{\mathcal I}\) is prime.  Moreover, the defining contraction equality for \(\mathcal I_{\mathrm{obs}}\) gives
\[
\overline{\mathcal I}\cap A=(0).
\]
Thus \(\overline{\mathcal I}\) is disjoint from \(A\setminus\{0\}\), and localizing at this multiplicative set produces a proper prime ideal
\[
\mathcal I_{\leq3}K_{\mathrm{obs}}[\beta]
\subsetneq K_{\mathrm{obs}}[\beta].
\]
The already proved relation and the nonzero class of \(W_2^{\mathrm{el}}\) imply
\[
\left\langle
\beta-\frac{W_3^{\mathrm{el}}}{W_2^{\mathrm{el}}}
\right\rangle
\subseteq
\mathcal I_{\leq3}K_{\mathrm{obs}}[\beta].
\]
The ideal on the left is maximal, since evaluation at \(W_3^{\mathrm{el}}/W_2^{\mathrm{el}}\in K_{\mathrm{obs}}\) gives
\[
K_{\mathrm{obs}}[\beta]
\bigg/
\left\langle
\beta-\frac{W_3^{\mathrm{el}}}{W_2^{\mathrm{el}}}
\right\rangle
\cong K_{\mathrm{obs}}.
\]
A proper ideal containing a maximal ideal equals that maximal ideal.  Therefore
\[
\mathcal I_{\leq3}K_{\mathrm{obs}}[\beta]
=
\left\langle
\beta-\frac{W_3^{\mathrm{el}}}{W_2^{\mathrm{el}}}
\right\rangle.
\]
Its quotient is the field \(K_{\mathrm{obs}}\), of vector-space dimension one over itself; equivalently, the generic scheme-theoretic \(\beta\)-fiber has degree one.  At any parameterized image point for which \(W_2^{\mathrm{el}}\ne0\), the identity \(R_{\mathrm{el}}=0\) can be divided by \(W_2^{\mathrm{el}}\), giving the pointwise formula
\[
\beta=\frac{W_3^{\mathrm{el}}}{W_2^{\mathrm{el}}}.
\]
Finally, every retained eliminant in the extended full ideal is a multiple of this linear generator.  An extra root of one such multiple can arise only from its multiplier and is not an additional common root of the full localized ideal.  The division and localization arguments do not apply on \(W_2^{\mathrm{el}}=0\); hence no classification of that exceptional divisor, and no identification assertion for the entire unrestricted real parameter class, follows from this proof.
\end{proof}
\par\smallskip\noindent \textit{Justification.} The theorem turns an independently derived cubic contraction into a precisely testable statement about the exact released full ideal and its observable localization. \textit{Gap.} Tramontano et al. (2025), Theorem 3.6, reports generic nonidentification through order three; the unreported object tested here is the localized full-ideal relation \(\beta W_2^{\mathrm{el}}-W_3^{\mathrm{el}}=0\) under the exact published parameterization. \textit{Consumer.} The result determines whether the one-proxy implementation may use a generic cubic effect chart on \(W_2^{\mathrm{el}}\ne0\), subject to the pinned replication gate.

\begin{theorem}[thm:covariance-overlap]\label{thm:covariance-overlap}
Let \(\pi_f=(1,2,3)\) and \(\pi_r=(3,2,1)\), and fix \(u=v=3/10\). For \(\pi_f\), take \(\gamma=(7/10,2/5,43027/61000)\) and \(s^2=(13/10,4/5,341693977/372100000)\); reverse the coordinates for \(\pi_r\). Both parameter points are interior to their positive boxes, \[\Psi_{\pi_f}(3/10,3/10,\gamma,s)=\Psi_{\pi_r}(3/10,3/10,\gamma^{\mathrm{rev}},s^{\mathrm{rev}})=\Sigma_0,\] and both six-by-six Jacobians of \((\gamma,s^2)\mapsto\Psi_\pi(3/10,3/10,\gamma,s)\) have nonzero determinant at these points. Therefore \(\Sigma_0\in\operatorname{int}(\mathcal S_{\pi_f}\cap\mathcal S_{\pi_r})\), so covariance alone fails to identify the order on a nonempty open set.
\end{theorem}
\begin{proof}
Put \(d_i=s_i^2\) and \(D_s=\operatorname{diag}(s)\).  By the structural equation in \(\mathrm{ass{:}structural\text{-}equation}\),
\[
 \operatorname{Cov}(X)
 =A_\pi\operatorname{Cov}(\gamma H+D_s\xi)A_\pi^{\mathsf T}.
\]
Mutual independence in \(\mathrm{ass{:}source\text{-}independence}\) makes all cross-covariances among \(H,\xi_1,\xi_2,\xi_3\) vanish, and \(\mathrm{ass{:}source\text{-}unit\text{-}variance}\) makes their four variances equal to one.  Hence
\[
 \operatorname{Cov}(\gamma H+D_s\xi)
 =\gamma\gamma^{\mathsf T}+D_s I_3D_s
 =\gamma\gamma^{\mathsf T}+\operatorname{diag}(d),
\]
and therefore
\[
 \operatorname{Cov}(X)
 =A_\pi\bigl(\gamma\gamma^{\mathsf T}+\operatorname{diag}(d)\bigr)A_\pi^{\mathsf T}
 =\Psi_\pi(u,v,\gamma,s).
 \tag{1}
\]

We first verify feasibility of the displayed forward parameter point.  Both coefficients are strictly interior because
\[
 \frac{3}{10}-\frac14=\frac1{20}>0,
 \qquad
 \frac34-\frac{3}{10}=\frac9{20}>0.
\]
For the loading coordinates, the only non-immediate endpoint checks are
\[
 \frac{43027}{61000}-\frac14=\frac{27777}{61000}>0,
 \qquad
 \frac34-\frac{43027}{61000}=\frac{2723}{61000}>0;
\]
also \(7/10,2/5\in(1/4,3/4)\).  Since \(s\) is the coordinatewise positive square root of \(d\), the residual-scale restriction is equivalent to \(d_i\in(9/16,25/16)\).  The six relevant strict differences are
\[
 \frac{13}{10}-\frac9{16}=\frac{59}{80},\quad
 \frac{25}{16}-\frac{13}{10}=\frac{21}{80},\quad
 \frac45-\frac9{16}=\frac{19}{80},\quad
 \frac{25}{16}-\frac45=\frac{61}{80},
\]
and
\[
 \frac{341693977}{372100000}-\frac9{16}
 =\frac{132387727}{372100000}>0,
 \qquad
 \frac{25}{16}-\frac{341693977}{372100000}
 =\frac{239712273}{372100000}>0.
\]
Thus the forward point satisfies the strict versions of \(\mathrm{ass{:}positive\text{-}mechanism\text{-}box}\) and \(\mathrm{ass{:}residual\text{-}scale\text{-}box}\).  Coordinate reversal preserves these intervals, so the reverse point is strict as well.

For \(\pi_f=(1,2,3)\), nilpotence of \(B_{\pi_f}\) gives
\[
 A_{\pi_f}=I_3+B_{\pi_f}+B_{\pi_f}^2
 =\begin{pmatrix}
 1&0&0\\[2pt]
 3/10&1&0\\[2pt]
 9/100&3/10&1
 \end{pmatrix}.
\]
Writing \(\lambda=A_{\pi_f}\gamma\), exact multiplication yields
\[
 \lambda
 =\begin{pmatrix}
 7/10\\[2pt]
 (3/10)(7/10)+2/5\\[2pt]
 (9/100)(7/10)+(3/10)(2/5)+43027/61000
 \end{pmatrix}
 =\begin{pmatrix}7/10\\[2pt]61/100\\[2pt]5419/6100\end{pmatrix}.
\]
Equation \((1)\) can consequently be written as
\[
 \Psi_{\pi_f}=\lambda\lambda^{\mathsf T}
 +A_{\pi_f}\operatorname{diag}(d)A_{\pi_f}^{\mathsf T}.
\]
Its six distinct entries are, by direct substitution,
\[
 \begin{aligned}
 (\Psi_{\pi_f})_{11}
 &=\left(\frac7{10}\right)^2+\frac{13}{10}=\frac{179}{100},\\
 (\Psi_{\pi_f})_{12}
 &=\frac7{10}\frac{61}{100}+\frac3{10}\frac{13}{10}=\frac{817}{1000},\\
 (\Psi_{\pi_f})_{13}
 &=\frac7{10}\frac{5419}{6100}+\frac9{100}\frac{13}{10}=\frac{4507}{6100},\\
 (\Psi_{\pi_f})_{22}
 &=\left(\frac{61}{100}\right)^2+\left(\frac3{10}\right)^2\frac{13}{10}+\frac45
 =\frac{12891}{10000},\\
 (\Psi_{\pi_f})_{23}
 &=\frac{61}{100}\frac{5419}{6100}
 +\left(\frac3{10}\right)^3\frac{13}{10}
 +\frac3{10}\frac45
 =\frac{817}{1000},\\
 (\Psi_{\pi_f})_{33}
 &=\left(\frac{5419}{6100}\right)^2
 +\left(\frac9{100}\right)^2\frac{13}{10}
 +\left(\frac3{10}\right)^2\frac45
 +\frac{341693977}{372100000}\\
 &=\frac{666059000}{372100000}=\frac{179}{100}.
 \end{aligned}
\]
These identities prove
\[
 \Psi_{\pi_f}(3/10,3/10,\gamma,s)=\Sigma_0.
 \tag{2}
\]
Moreover, \(\gamma\gamma^{\mathsf T}+\operatorname{diag}(d)\) is positive definite because every \(d_i>0\), and \(A_{\pi_f}\) is invertible.  Thus \((2)\) also verifies directly that \(\Sigma_0\in\operatorname{Sym}_{++}^3\).

Let
\[
 R=\begin{pmatrix}0&0&1\\0&1&0\\1&0&0\end{pmatrix}
\]
be coordinate reversal.  From the definition of \(B_\pi\),
\[
 B_{\pi_r}(3/10,3/10)=RB_{\pi_f}(3/10,3/10)R,
 \qquad
 A_{\pi_r}=RA_{\pi_f}R.
\]
The reversed nuisance parameters obey
\[
 \gamma^{\mathrm{rev}}=R\gamma,
 \qquad
 \operatorname{diag}((s^{\mathrm{rev}})^2)
 =R\operatorname{diag}(d)R.
\]
It follows by associativity and \(R^{\mathsf T}=R=R^{-1}\) that
\[
 \begin{aligned}
 \Psi_{\pi_r}(3/10,3/10,\gamma^{\mathrm{rev}},s^{\mathrm{rev}})
 &=R\Psi_{\pi_f}(3/10,3/10,\gamma,s)R\\
 &=R\Sigma_0R
 =\Sigma_0,
 \end{aligned}
 \tag{3}
\]
where the last equality follows from the equalities \((\Sigma_0)_{11}=(\Sigma_0)_{33}\) and \((\Sigma_0)_{12}=(\Sigma_0)_{23}\) visible in \(\mathrm{def{:}covariance\text{-}witness}\).

It remains to prove that the common covariance is an interior, rather than isolated, intersection point.  For a loading \(\eta=(\eta_a,\eta_b,\eta_c)\) in chain order, define
\[
 F(\eta,d)=\eta\eta^{\mathsf T}+\operatorname{diag}(d).
\]
In output coordinates \((F_{aa},F_{bb},F_{cc},F_{ab},F_{ac},F_{bc})\) and input coordinates \((\eta_a,\eta_b,\eta_c,d_a,d_b,d_c)\), its Jacobian is
\[
 J_F=
 \begin{pmatrix}
 2\eta_a&0&0&1&0&0\\
 0&2\eta_b&0&0&1&0\\
 0&0&2\eta_c&0&0&1\\
 \eta_b&\eta_a&0&0&0&0\\
 \eta_c&0&\eta_a&0&0&0\\
 0&\eta_c&\eta_b&0&0&0
 \end{pmatrix},
 \qquad
 \det J_F=2\eta_a\eta_b\eta_c.
 \tag{4}
\]
The derivative of \((\gamma,d)\mapsto\Psi_\pi(3/10,3/10,\gamma,\sqrt d)\) is the composition of \(J_F\) with the congruence map \(Q\mapsto A_\pi Q A_\pi^{\mathsf T}\).  The latter linear map is invertible, with inverse \(Q\mapsto A_\pi^{-1}Q A_\pi^{-\mathsf T}\).  At the forward point, equation \((4)\) gives
\[
 2\gamma_1\gamma_2\gamma_3
 =2\frac7{10}\frac25\frac{43027}{61000}
 =\frac{301189}{762500}>0.
\]
At the reverse point the three factors are merely permuted.  Consequently both six-by-six Jacobian determinants in the theorem are nonzero.

For completeness, an explicit inverse proves the required open-image assertion.  Fix either \(\pi\in\{\pi_f,\pi_r\}\), keep \(u=v=3/10\), and, for a candidate covariance \(S\), set
\[
 \Omega_\pi(S)=A_\pi^{-1}S A_\pi^{-\mathsf T}.
\]
Whenever the three off-diagonal entries of \(\Omega=\Omega_\pi(S)\), written in chain order, are positive, define
\[
 \eta_a(\Omega)=\sqrt{\frac{\Omega_{ab}\Omega_{ac}}{\Omega_{bc}}},
 \qquad
 \eta_b(\Omega)=\frac{\Omega_{ab}}{\eta_a(\Omega)},
 \qquad
 \eta_c(\Omega)=\frac{\Omega_{ac}}{\eta_a(\Omega)},
 \qquad
 d_i(\Omega)=\Omega_{ii}-\eta_i(\Omega)^2.
 \tag{5}
\]
The equations in \((5)\) give, one by one,
\[
 \eta_a\eta_b=\Omega_{ab},\qquad
 \eta_a\eta_c=\Omega_{ac},\qquad
 \eta_b\eta_c
 =\frac{\Omega_{ab}\Omega_{ac}}{\eta_a^2}
 =\Omega_{bc},
\]
and the definition of \(d_i\) gives \(\eta_i^2+d_i=\Omega_{ii}\).  Hence
\[
 F(\eta(\Omega),d(\Omega))=\Omega.
 \tag{6}
\]
At \(S=\Sigma_0\), formulas \((5)\) recover exactly the applicable displayed loading and squared-scale vectors, because all loading coordinates are positive and \(\Omega_{ij}=\gamma_i\gamma_j\) for \(i\ne j\).  The maps in \((5)\) are continuous wherever those three off-diagonal entries are positive.  The strict inequalities proved above therefore give an open neighborhood \(V_\pi\) of \(\Sigma_0\) on which every recovered \(\eta_i\) lies in \((1/4,3/4)\) and every recovered \(d_i\) lies in \((9/16,25/16)\).  For \(S\in V_\pi\), take the coordinatewise positive square root \(s_i=\sqrt{d_i}\).  Equations \((5)\)--\((6)\) then imply
\[
 S=A_\pi F(\eta,d)A_\pi^{\mathsf T}
 =\Psi_\pi(3/10,3/10,\eta,s),
\]
so \(V_\pi\subseteq\mathcal S_\pi\).  Thus
\[
 \Sigma_0\in V_{\pi_f}\cap V_{\pi_r}
 \subseteq\mathcal S_{\pi_f}\cap\mathcal S_{\pi_r},
\]
and \(V_{\pi_f}\cap V_{\pi_r}\) is an open neighborhood of \(\Sigma_0\).  Therefore \(\Sigma_0\in\operatorname{int}(\mathcal S_{\pi_f}\cap\mathcal S_{\pi_r})\).  Since \(\pi_f\ne\pi_r\) and every covariance in this nonempty open neighborhood is compatible with both labeled chains, no functional of covariance alone can identify the chain order there.
\end{proof}
\par\smallskip\noindent \textit{Justification.} An open overlap proves that the cubic tensor, rather than covariance information already implicit in the model, performs the order identification. \textit{Gap.} Ghassami et al. (2020) characterize distribution equivalence abstractly, while Foygel Barber et al. (2022) treat known-graph covariance recovery; neither gives this pairwise positive-chain image overlap. \textit{Consumer.} The overlap supplies a rigorous ablation benchmark for any implementation claiming that the chain label is learned from skewness.

\begin{theorem}[thm:honest-inversion]\label{thm:honest-inversion}
For every \(n\ge1\) and \(0<\alpha<1\), \[\inf_{P\in\mathcal P_+}\mathbb P_P\!\left(B_\pi(u,v)\in\mathsf{CS}_n(\alpha)\right)\ge1-\alpha.\] If \(2\varepsilon_n(\alpha)<\delta_0\), then on \(\{\|\widehat C-C\|_\infty\le\varepsilon_n(\alpha)\}\), every element of \(\mathsf{CS}_n(\alpha)\) has the true label and \[\operatorname{diam}_\infty(\mathsf{CS}_n(\alpha))\le2L_\star\varepsilon_n(\alpha)=O(n^{-1/2}).\] These conclusions hold for every rank of \(\operatorname{Var}_P(\operatorname{vech}(X^{\otimes3}))\).
\end{theorem}
\begin{proof}
Fix \(P\in\mathcal P_+\), \(n\geq1\), and \(\alpha\in(0,1)\), and choose a representation of \(P\) furnished by \(\texttt{def:positive-model}\), with causal order \(\pi=(a,b,c)\).  We first obtain a support bound uniform over that representation.  Since the two-edge chain is acyclic, \(B_\pi^3=0\) and hence \(A_\pi=I_3+B_\pi+B_\pi^2\).  Substituting this identity into \(\texttt{ass:structural-equation}\), and using \(x=\gamma_a\), \(y=\gamma_b\), and \(z=\gamma_c\), gives
\[
 X_a=xH+s_a\xi_a,
 \qquad
 X_b=uX_a+yH+s_b\xi_b,
 \qquad
 X_c=vX_b+zH+s_c\xi_c.
\]
By \(\texttt{ass:source-support}\), every source on the right-hand side has absolute value at most \(2\) almost surely.  By \(\texttt{ass:positive-mechanism-box}\) and \(\texttt{ass:residual-scale-box}\),
\[
 |X_a|
 \leq 2(x+s_a)
 \leq2\left(\frac34+\frac54\right)=4,
\]
\[
 |X_b|
 \leq u|X_a|+2(y+s_b)
 \leq\frac34\,4+2\left(\frac34+\frac54\right)=7,
\]
and
\[
 |X_c|
 \leq v|X_b|+2(z+s_c)
 \leq\frac34\,7+2\left(\frac34+\frac54\right)
 =\frac{37}{4}=R_X.
\]
Thus
\[
 \max_{i\in\mathcal V}|X_i|\leq R_X
 \quad\text{almost surely}.                                      \tag{1}
\]
By \(\texttt{ass:source-centering}\) and the linear structural equation, \(\mathbb E_P[X]=0\).  The defining third-cumulant identity therefore reduces, for all \(i,j,k\in\mathcal V\), to
\[
 C_{ijk}=\operatorname{Cum}(X_i,X_j,X_k)
 =\mathbb E_P[X_iX_jX_k].                                         \tag{2}
\]

Fix one of the ten unordered triples \(1\leq i\leq j\leq k\leq3\), and define
\[
 Z_\ell=X_i^{(\ell)}X_j^{(\ell)}X_k^{(\ell)},
 \qquad 1\leq\ell\leq n.
\]
The variables \(Z_1,\ldots,Z_n\) are independent by \(\texttt{ass:iid-sampling}\), and (1) gives
\[
 -R_X^3\leq Z_\ell\leq R_X^3
 \quad\text{almost surely}.
\]
Apply \(\texttt{lem:hoeffding-bounded-sum}\) to \(Z_1,\ldots,Z_n\) and then to \(-Z_1,\ldots,-Z_n\), in each case with deviation \(n\varepsilon\).  Equations (2) and \(\texttt{def:empirical-cubic}\) then imply, for every \(\varepsilon>0\),
\[
 \begin{aligned}
 \mathbb P_P\!\left(
   |\widehat C_{ijk}-C_{ijk}|>\varepsilon
 \right)
 &\leq
 2\exp\!\left(
 -\frac{2n^2\varepsilon^2}{n(2R_X^3)^2}
 \right)\\
 &=2\exp\!\left(-\frac{n\varepsilon^2}{2R_X^6}\right).
 \end{aligned}                                                     \tag{3}
\]
A union bound over the ten coordinates of a symmetric cubic tensor yields
\[
 \mathbb P_P\!\left(\|\widehat C-C\|_\infty>\varepsilon\right)
 \leq20\exp\!\left(-\frac{n\varepsilon^2}{2R_X^6}\right).        \tag{4}
\]
At
\[
 \varepsilon=\varepsilon_n(\alpha)
 =R_X^3\sqrt{\frac{2\log(20/\alpha)}{n}},
\]
the right-hand side of (4) equals \(\alpha\).  Put
\[
 \mathcal E_n
 =\{\|\widehat C-C\|_\infty\leq\varepsilon_n(\alpha)\}.
\]
The generating parameter \(\theta\) belongs to \(\Theta_+\) by \(\texttt{def:positive-model}\) and \(\texttt{def:positive-parameter-set}\), while \(\texttt{lem:cubic-factorization}\) gives \(\Phi_+(\theta)=C\).  Hence \(\texttt{def:confidence-set}\) gives, on \(\mathcal E_n\),
\[
 B_\pi(u,v)\in\mathsf{CS}_n(\alpha).                              \tag{5}
\]
Since the bound (4) is uniform in \(P\in\mathcal P_+\), (5) proves
\[
 \inf_{P\in\mathcal P_+}
 \mathbb P_P\!\left(B_\pi(u,v)\in\mathsf{CS}_n(\alpha)\right)
 \geq
 \inf_{P\in\mathcal P_+}\mathbb P_P(\mathcal E_n)
 \geq1-\alpha.                                                     \tag{6}
\]

We next establish the deterministic separation and inverse-Lipschitz bounds needed for the diameter claim.  Let
\[
 C^{(r)}=\Phi_+(\theta^{(r)}),
 \qquad \theta^{(r)}\in\Theta_+,
 \qquad r\in\{0,1\},
\]
and write
\[
 \eta=\|C^{(0)}-C^{(1)}\|_\infty.
\]
For each \(r\), apply \(\texttt{prop:source-feasibility}\) separately to the four coordinates of the skewness vector in \(\theta^{(r)}\), and take the product of the resulting four independent laws.  Together with the remaining coordinates of \(\theta^{(r)}\), these sources and \(\texttt{ass:structural-equation}\) produce a law \(P_r\in\mathcal P_+\).  Lemma \(\texttt{lem:cubic-factorization}\) identifies its cubic tensor with \(C^{(r)}\).  Applying (1)--(2) to \(P_r\) shows that every entry of \(C^{(r)}\), and therefore every entry of each of its slices, has absolute value at most
\[
 K_0=R_X^3.                                                        \tag{7}
\]

We use only elementary polynomial perturbation bounds.  If \(|q_j|,|q'_j|\leq K_0\) and \(|q_j-q'_j|\leq\eta\) for \(j=1,2,3\), telescoping one factor at a time gives
\[
 |q_1q_2q_3-q'_1q'_2q'_3|
 \leq3K_0^2\eta.                                                  \tag{8}
\]
A \(3\)-by-\(3\) determinant is a signed sum of six triple products, so (8) gives, for corresponding slices \(M,M'\),
\[
 |\det(M)-\det(M')|\leq18K_0^2\eta.                              \tag{9}
\]
Each adjugate entry is a signed sum of two products of two entries.  The analogous two-factor telescoping identity and (7) give
\[
 \|\operatorname{adj}(M)-\operatorname{adj}(M')\|_\infty
 \leq4K_0\eta.                                                    \tag{10}
\]

Suppose first that the root labels of \(C^{(0)}\) and \(C^{(1)}\) differ, and let \(a_0\) be the root label of \(C^{(0)}\).  By \(\texttt{thm:global-cubic-inverse}\),
\[
 \det(M_{a_0}^{(0)})=0,
 \qquad
 |\det(M_{a_0}^{(1)})|\geq\min(d_b,d_c),
\]
because \(a_0\) is a nonroot label for \(C^{(1)}\).  Inequality (9) therefore implies
\[
 \eta\geq\frac{\min(d_b,d_c)}{18K_0^2}.                           \tag{11}
\]
Suppose instead that the roots agree, with common value \(a\), but the middle labels differ.  Let \(j\) be the middle label for \(C^{(0)}\); because there are three labels, \(j\) is the terminal label for \(C^{(1)}\).  Lemma \(\texttt{lem:adjugate-sign-margin}\) gives
\[
 \bigl(\operatorname{adj}(M_a^{(0)})\bigr)_{aj}\leq-m,
 \qquad
 \bigl(\operatorname{adj}(M_a^{(1)})\bigr)_{aj}\geq m.
\]
Combining the resulting difference of at least \(2m\) with (10) yields
\[
 \eta\geq\frac{m}{2K_0}.                                         \tag{12}
\]
By (11), (12), and the definition of \(\delta_0\),
\[
 \|C^{(0)}-C^{(1)}\|_\infty<\delta_0
 \quad\Longrightarrow\quad
 \pi^{(0)}=\pi^{(1)}.                                             \tag{13}
\]

It remains to control the recovered coefficients when the common label is \(\pi=(a,b,c)\).  For \(r\in\{0,1\}\), set
\[
 N^{(r)}=\operatorname{adj}(M_a^{(r)})e_c,
 \qquad
 W^{(r)}=M_b^{(r)}N^{(r)},
\]
and define
\[
 D^{(r)}=N_b^{(r)}W_b^{(r)}+N_c^{(r)}W_c^{(r)},
 \qquad
 A^{(r)}=-N_a^{(r)}W_b^{(r)}.
\]
The two-term cofactor formula, (7), and (10) imply
\[
 \|N^{(r)}\|_\infty\leq2K_0^2,
 \qquad
 \|N^{(0)}-N^{(1)}\|_\infty\leq4K_0\eta.                       \tag{14}
\]
Expanding each of the three terms in the row--column product \(M_bN\), and adding and subtracting the corresponding cross term, gives from (7) and (14)
\[
 \begin{aligned}
 \|W^{(0)}-W^{(1)}\|_\infty
 &\leq3\left(2K_0^2\eta+4K_0^2\eta\right)
 =18K_0^2\eta,\\
 \|W^{(r)}\|_\infty
 &\leq3K_0(2K_0^2)=6K_0^3.
 \end{aligned}                                                     \tag{15}
\]
Using (14)--(15) once for the single product defining \(A^{(r)}\), and twice for the two products defining \(D^{(r)}\), yields
\[
 \begin{aligned}
 |A^{(0)}-A^{(1)}|
 &\leq(4K_0\eta)(6K_0^3)+(2K_0^2)(18K_0^2\eta)
 =60K_0^4\eta,\\
 |D^{(0)}-D^{(1)}|
 &\leq120K_0^4\eta.
 \end{aligned}                                                     \tag{16}
\]

The denominator bounds and recovery identities in \(\texttt{thm:global-cubic-inverse}\) give
\[
 |W_b^{(r)}|\geq d_W,
 \qquad |D^{(r)}|\geq d_D,
 \qquad
 v^{(r)}=\frac{W_c^{(r)}}{W_b^{(r)}},
 \qquad
 u^{(r)}=\frac{A^{(r)}}{D^{(r)}}.                                  \tag{17}
\]
For nonzero \(s,s'\), the exact identity
\[
 \frac{r}{s}-\frac{r'}{s'}
 =\frac{r-r'-(r'/s')(s-s')}{s}
\]
implies
\[
 \left|\frac{r}{s}-\frac{r'}{s'}\right|
 \leq\frac{|r-r'|+|r'/s'|\,|s-s'|}{|s|}.                         \tag{18}
\]
Apply (18) first to the ratio for \(v\) in (17), using (15), and then to the ratio for \(u\), using (16).  Since \(\theta^{(r)}\in\Theta_+\) implies \(|u^{(r)}|,|v^{(r)}|\leq3/4\), we obtain
\[
 |v^{(0)}-v^{(1)}|
 \leq
 \frac{18K_0^2+(3/4)18K_0^2}{d_W}\eta
 =\frac{63K_0^2}{2d_W}\eta,                                      \tag{19}
\]
and
\[
 |u^{(0)}-u^{(1)}|
 \leq
 \frac{60K_0^4+(3/4)120K_0^4}{d_D}\eta
 =\frac{150K_0^4}{d_D}\eta.                                      \tag{20}
\]
On a common labeled chart, the only nonzero entries of
\(B_\pi(u^{(0)},v^{(0)})-B_\pi(u^{(1)},v^{(1)})\) are the two coefficient differences.  Hence (13), (19), (20), and the definition of \(L_\star\) prove
\[
 \|C^{(0)}-C^{(1)}\|_\infty<\delta_0
 \quad\Longrightarrow\quad
 \begin{cases}
 \pi^{(0)}=\pi^{(1)},\\[2pt]
 \|B_{\pi^{(0)}}(u^{(0)},v^{(0)})
      -B_{\pi^{(1)}}(u^{(1)},v^{(1)})\|_\infty
 \leq L_\star\|C^{(0)}-C^{(1)}\|_\infty.
 \end{cases}                                                       \tag{21}
\]

Assume now that \(2\varepsilon_n(\alpha)<\delta_0\) and that \(\mathcal E_n\) occurs.  If \(\theta'\in\Theta_+\) is retained in the defining feasibility problem for \(\mathsf{CS}_n(\alpha)\), then the triangle inequality and \(\texttt{def:confidence-set}\) give
\[
 \|\Phi_+(\theta')-C\|_\infty
 \leq
 \|\Phi_+(\theta')-\widehat C\|_\infty
 +\|\widehat C-C\|_\infty
 \leq2\varepsilon_n(\alpha)<\delta_0.
\]
Apply (21) to \(\theta'\) and the true generating parameter.  Every retained parameter consequently has the true label.  Moreover, for any two retained parameters \(\theta^{(0)},\theta^{(1)}\),
\[
 \|\Phi_+(\theta^{(0)})-\Phi_+(\theta^{(1)})\|_\infty
 \leq2\varepsilon_n(\alpha)<\delta_0.
\]
Applying (21) once more and taking the supremum over all pairs of retained matrices proves
\[
 \operatorname{diam}_\infty(\mathsf{CS}_n(\alpha))
 \leq2L_\star\varepsilon_n(\alpha).                               \tag{22}
\]
For fixed \(\alpha\), the definition of \(\varepsilon_n(\alpha)\) makes the right-hand side of (22) \(O(n^{-1/2})\).

Finally, the probability argument (1)--(6) uses only bounded source support, source centering, and independence across observations; the deterministic argument (7)--(22) uses polynomial perturbation inequalities and the uniform scalar denominator bounds in \(\texttt{thm:global-cubic-inverse}\).  At no step is \(\operatorname{Var}_P(\operatorname{vech}(X^{\otimes3}))\) inverted or required to have any rank.  Thus (6) and (22), including the label conclusion, remain valid for every rank of that covariance matrix.
\end{proof}
\par\smallskip\noindent \textit{Justification.} The set gives honest uncertainty quantification across discrete order selection and rational denominators without a Wald approximation. \textit{Gap.} Tramontano et al. (2025) provide point estimators but no coverage theorem; Babii (2020) and Andrews and Guggenberger (2019) do not treat a finite union of exact LvLiNGAM charts. \textit{Consumer.} CEId-from-Moments can return a certified effect set and an order-ambiguity warning even when empirical moment covariance is singular.

\begin{proposition}[prop:fourth-order-target-compatibility]\label{prop:fourth-order-target-compatibility}
As a direct target-compatibility corollary, on the fixed order \(\pi=(1,2,3)\), \(\operatorname{Rec}(C)=B_\pi(u,v)\) returns \(v=B_{\pi,32}\). For the same \(P\in\mathcal P_+\), the causal effect identified in Tramontano et al. (2025), Theorem 3.5, is defined as this same structural coordinate.
\end{proposition}
\begin{proof}
Fix \(P\in\mathcal P_+\) whose labeled order is \(\pi=(1,2,3)\), and let \(C=C_3(P)\).  By the global cubic inverse theorem,
\[
\operatorname{Rec}(C)=B_\pi(u,v).
\]
This equality is a proved identification result, not a definition of the causal coefficient.  From the definition of the structural coefficient matrix and the identities \(a=1\), \(b=2\), and \(c=3\),
\[
B_\pi(u,v)=u e_2e_1^{\mathsf T}+v e_3e_2^{\mathsf T}.
\]
Therefore, using \(e_i^{\mathsf T}e_j=\mathbf 1\{i=j\}\),
\[
\begin{aligned}
\bigl[\operatorname{Rec}(C)\bigr]_{32}
&=e_3^{\mathsf T}\operatorname{Rec}(C)e_2\\
&=e_3^{\mathsf T}B_\pi(u,v)e_2\\
&=u(e_3^{\mathsf T}e_2)(e_1^{\mathsf T}e_2)
  +v(e_3^{\mathsf T}e_3)(e_2^{\mathsf T}e_2)\\
&=v.
\end{aligned}
\]
Hence the cubic recovery map returns \(v=(B_\pi)_{32}\).

It remains only to match notation for the causal target.  By the published target-coordinate lemma, the target in Tramontano et al. (2025), Theorem 3.5, is the structural coefficient \(b_{Y,T}\) on the directed edge from the observed treatment \(T\) to the observed outcome \(Y\).  In the fixed published graph encoded by \(L_{\mathrm{pub}}\) and \(\Lambda_{\mathrm{pub}}\), the treatment and outcome are observed nodes \(2\) and \(3\), respectively, and the coefficient on \(2\to3\) is \(\beta\).  Under the order \(\pi=(1,2,3)\), that same edge coefficient is the \((3,2)\) entry of \(B_\pi\).  Consequently,
\[
b_{Y,T}=\beta=(B_\pi)_{32}=v
=\bigl[\operatorname{Rec}(C)\bigr]_{32}.
\]
Thus the published fourth-order result and the cubic inverse concern the same causal structural coordinate.  The argument asserts compatibility of their targets only; it does not identify their distinct observable recovery functionals with one another.
\end{proof}
\par\smallskip\noindent \textit{Justification.} This bookkeeping corollary records target compatibility and requires no independent fourth-order root solve. \textit{Gap.} The published fourth-order result and the cubic inverse use different information, so their common structural coordinate must be stated without treating equality of targets as a new identification theorem. \textit{Consumer.} The compatibility remark preserves the treatment-to-outcome interpretation when the cubic chart is used in the published workflow.

\begin{openendedquestion}[oeq:real-fiber-atlas]\label{oeq:real-fiber-atlas}
Construct \(\mathfrak A\) to complete the exceptional classification not covered by the principal real charts.  First, for the cubic map \((\pi,h,r,t,w,x,y,z,u,v)\mapsto C_{\pi}^{(3)}(h,r,t,w,x,y,z,u,v)\), partition
\[
\bigsqcup_{\pi\in\mathfrak S_3}\{\pi\}\times
\bigl(\mathbb R^9\setminus\Theta_{\pi}^{\circ}\bigr)
\]
into finitely many semialgebraic strata on which the real fiber cardinality and the set of compatible chain orders are certified and constant.  On every singleton-order stratum, give observable rational formulas for \(u\) and \(v\); on every non-singleton or positive-dimensional stratum, give explicit real parameter twins or a fiber parameterization.  Second, without conflating this six-order cubic problem with the fixed published graph, stratify the exceptional complement of the analytically certified localized chart \(W_2^{\mathrm{el}}\ne0\) for the observed-moment fibers underlying \(\widetilde\Phi_{\le3}\), recording all compatible values of \(\beta\).  On all intersections of slice-rank, kernel-coordinate, adjugate-denominator, and localization-denominator hypersurfaces, decide whether failure of the displayed \(\operatorname{Rec}\) branch is removable by an alternate chart or reflects genuine nonidentification in the relevant cubic or full second-and-third-moment map.
\end{openendedquestion}
\par\smallskip\noindent \textit{Justification.} This exposes the genuine unresolved content instead of equating one failed rational denominator with structural ambiguity. \textit{Gap.} Generic elimination in Tramontano et al. (2025) and Garcia-Puente et al. (2010) does not classify real fibers on exceptional intersections; Drton, Foygel, and Sullivant (2011) do not cover this higher-moment latent LiNGAM map. \textit{Consumer.} A completed atlas yields exact ambiguity warnings and alternate formulas for observations near every failure divisor of the cubic branch.

\begin{lemma}[lem:adjugate-sign-margin]\label{lem:adjugate-sign-margin}
For every \(P\in\mathcal P_+\) with \(\pi=(a,b,c)\), the root-slice adjugate satisfies
\[
\bigl(\operatorname{adj}(M_a)\bigr)_{ab}<0,
\qquad
\bigl(\operatorname{adj}(M_a)\bigr)_{ac}>0,
\]
and
\[
\min\!\left\{\left|\bigl(\operatorname{adj}(M_a)\bigr)_{ab}\right|,
\left|\bigl(\operatorname{adj}(M_a)\bigr)_{ac}\right|\right\}\ge m.
\]
\end{lemma}
\begin{proof}
Put \(K=hxr\).  The two-direction representation
\(M_a=hxgg^{\mathsf T}+rdd^{\mathsf T}\) from the cubic factorization implies, by direct expansion of its \(2\times2\) cofactors,
\[
\operatorname{adj}(M_a)=K(g\times d)(g\times d)^{\mathsf T}=Kkk^{\mathsf T}.
\]
Because \(k=(-uz,z+vy,-y)^{\mathsf T}\),
\[
\bigl(\operatorname{adj}(M_a)\bigr)_{ab}
=-Kuz(z+vy)<0,
\qquad
\bigl(\operatorname{adj}(M_a)\bigr)_{ac}
=Kuzy>0.
\]
The box assumptions give
\[
K=hxr\ge \frac12\cdot\frac14\cdot
\left(\frac34\right)^3\frac12
=\frac{27}{1024},
\]
and \(u,y,z\ge1/4\), while \(z+vy\ge z\ge1/4\).  Each of the preceding absolute values is therefore at least
\[
\frac{27}{1024}\left(\frac14\right)^3
=\frac{27}{65536}=m.
\]
\end{proof}

\begin{lemma}[lem:hoeffding-bounded-sum]\label{lem:hoeffding-bounded-sum}
Let \(J\ge1\), and let \(Z_1,\ldots,Z_J\) be independent real random variables such that \(A_j\le Z_j\le B_j\) almost surely.  Then, for every \(t>0\),
\[
\mathbb P\!\left(\sum_{j=1}^J\bigl(Z_j-\mathbb E[Z_j]\bigr)\ge t\right)
\le
\exp\!\left(-\frac{2t^2}{\sum_{j=1}^J(B_j-A_j)^2}\right).
\]
\end{lemma}

\begin{lemma}[lem:published-proxy-target-coordinate]\label{lem:published-proxy-target-coordinate}
For the one-proxy graph in Figure 3 of Tramontano et al. (2025), the causal target in their Theorem 3.5 is the structural coefficient on the directed edge from the observed treatment \(T\) to the observed outcome \(Y\), denoted there by \(b_{Y,T}\).
\end{lemma}

\begin{lemma}[lem:principal-chart-algebra]\label{lem:principal-chart-algebra}
Fix \(\pi=(a,b,c)\in\mathfrak S_3\) and arbitrary \((h,r,t,w,x,y,z,u,v)\in\mathbb R^9\).  Put \(p=ux+y\), \(q=vp+z\), and \(C=C_{\pi}^{(3)}(h,r,t,w,x,y,z,u,v)\).  In coordinates ordered as \((a,b,c)\), let
\[
G=(x,p,q)^{\mathsf T},\quad R=(1,u,uv)^{\mathsf T},\quad E=(0,1,v)^{\mathsf T},\quad F=(0,0,1)^{\mathsf T},
\quad k=G\times R=(-uz,z+vy,-y)^{\mathsf T},
\quad K=hxr.
\]
Then
\[
M_a=hxGG^{\mathsf T}+rRR^{\mathsf T},\quad
M_b=hpGG^{\mathsf T}+ruRR^{\mathsf T}+tEE^{\mathsf T},\quad
M_c=hqGG^{\mathsf T}+ruvRR^{\mathsf T}+tvEE^{\mathsf T}+wFF^{\mathsf T},
\]
\[
\det(M_b)=hrtupz^2,
\]
\[
\det(M_c)=hrtuv^2qz^2+hrwuvqy^2+htwvqx^2+rtwuv^2,
\]
and
\[
\operatorname{adj}(M_a)=Kkk^{\mathsf T},\qquad
N=\operatorname{adj}(M_a)e_c
=K(uyz,-y(z+vy),y^2)^{\mathsf T},
\]
\[
W=M_bN=-Ktyz(0,1,v)^{\mathsf T},\qquad
W_b=-Ktyz,qquad W_c=-Ktvyz,qquad
D=K^2ty^2z^2,qquad -N_aW_b=K^2tuy^2z^2.
\]
\end{lemma}
\begin{proof}
By the definition of \(C_{\pi}^{(3)}\), contraction in its first slot gives, for every \(i\in\mathcal V\),
\[
M_i=hG_iGG^{\mathsf T}+rR_iRR^{\mathsf T}+tE_iEE^{\mathsf T}+wF_iFF^{\mathsf T}.
\]
Substitution of the \((a,b,c)\)-coordinates of \(G,R,E,F\) yields the three displayed slice formulas.  Moreover,
\[
G\times R
=\bigl(puv-qu,\ q-xuv,\ xu-p\bigr)^{\mathsf T}
=(-uz,z+vy,-y)^{\mathsf T}=k,
\]
where the second equality uses \(p=ux+y\) and \(q=vp+z\).

For three vectors \(v_1,v_2,v_3\in\mathbb R^3\) and scalars \(\lambda_1,\lambda_2,\lambda_3\), direct multiplication of
\(V\operatorname{diag}(\lambda_1,\lambda_2,\lambda_3)V^{\mathsf T}\), with columns of \(V\) equal to the \(v_j\), gives
\[
\det\!\left(\sum_{j=1}^3\lambda_jv_jv_j^{\mathsf T}\right)
=\lambda_1\lambda_2\lambda_3\det(v_1,v_2,v_3)^2.
\]
Here
\[
\det(G,R,E)=z,
\]
so the formula with weights \(hp,ru,t\) proves
\[
\det(M_b)=hp\,ru\,t\,z^2=hrtupz^2.
\]
Expanding the determinant of a sum of four rank-one matrices by multilinearity, all terms with a repeated selected vector vanish, and the four surviving three-vector terms are
\[
\det(M_c)
=hq\,ruv\,tv\det(G,R,E)^2
+hq\,ruv\,w\det(G,R,F)^2
+hq\,tv\,w\det(G,E,F)^2
+ruv\,tv\,w\det(R,E,F)^2.
\]
The four determinants are respectively \(z,-y,x,1\).  Substitution proves
\[
\det(M_c)=hrtuv^2qz^2+hrwuvqy^2+htwvqx^2+rtwuv^2.
\]

Since \(M_a=hxGG^{\mathsf T}+rRR^{\mathsf T}\), direct cofactor expansion gives
\[
\operatorname{adj}(M_a)=hxr(G\times R)(G\times R)^{\mathsf T}=Kkk^{\mathsf T}.
\]
Because \(k_c=-y\), multiplication by \(e_c\) gives
\[
N=Kk k_c=K(uyz,-y(z+vy),y^2)^{\mathsf T}.
\]
The cross-product identity implies \(G^{\mathsf T}N=R^{\mathsf T}N=0\), while direct substitution gives
\[
E^{\mathsf T}N=N_b+vN_c
=-Ky(z+vy)+Kvy^2=-Kyz.
\]
Consequently the formula for \(M_b\) yields
\[
W=M_bN=tE(E^{\mathsf T}N)=-Ktyz(0,1,v)^{\mathsf T}.
\]
Thus \(W_b=-Ktyz\) and \(W_c=-Ktvyz\).  Finally,
\[
D=N_bW_b+N_cW_c
=K^2ty^2z(z+vy)-K^2tvy^3z
=K^2ty^2z^2,
\]
and
\[
-N_aW_b=K^2tuy^2z^2.
\]
This proves every claimed polynomial identity without a sign or nonvanishing assumption.
\end{proof}

\begin{theorem}[thm:principal-real-cubic-inverse]\label{thm:principal-real-cubic-inverse}
For every \(\pi=(a,b,c)\in\mathfrak S_3\) and every \((h,r,t,w,x,y,z,u,v)\in\Theta_{\pi}^{\circ}\), let
\(C=C_{\pi}^{(3)}(h,r,t,w,x,y,z,u,v)\).  Then \(M_a\) has rank two, \(M_b\) and \(M_c\) are nonsingular, and hence \(a\) is the unique rank-two slice label.  Every nonzero \(\widetilde k\in\ker(M_a)\) has nonzero coordinates, and \(b\) is the unique label different from \(a\) whose coordinate has sign opposite to \(\widetilde k_a\).  For the resulting \(c\), the quantities in the principal recovery rule satisfy \(W_bD\ne0\) and
\[
v=\frac{W_c}{W_b},
\qquad
u=-\frac{N_aW_b}{D},
\qquad
\operatorname{Rec}(C)=B_{\pi}(u,v).
\]
If the same tensor has another principal-chart representation
\(C=C_{\pi'}^{(3)}(h',r',t',w',x',y',z',u',v')\) with parameters in \(\Theta_{\pi'}^{\circ}\), then \(\pi'=\pi\), \(u'=u\), and \(v'=v\).  Thus the unknown labeled order and both edge coefficients are globally identified on \(\mathcal C^{\circ}\), and the recovery uses constant arithmetic cost.
\end{theorem}
\begin{proof}
Fix a principal-chart representation.  By its defining product condition,
\[
h\ne0,\quad r\ne0,\quad t\ne0,\quad x\ne0,\quad y\ne0,\quad z\ne0,\quad u\ne0,\quad p\ne0,
\quad \det(M_c)\ne0.
\]
In particular, \(K=hxr\ne0\).  The two strict sign conditions also imply \(z+vy\ne0\).  Hence every coordinate of
\[
k=(-uz,z+vy,-y)^{\mathsf T}
\]
is nonzero.  Because \(k=G\times R\), one has \(G^{\mathsf T}k=R^{\mathsf T}k=0\).  The slice formula in the principal algebra lemma therefore gives
\[
M_ak=hxG(G^{\mathsf T}k)+rR(R^{\mathsf T}k)=0,
\]
while the same lemma gives \(\operatorname{adj}(M_a)=Kkk^{\mathsf T}\ne0\).  A singular \(3\times3\) matrix with nonzero adjugate has rank exactly two: singularity gives rank at most two, while a nonzero cofactor is a nonzero \(2\times2\) minor and gives rank at least two.  Therefore
\[
\operatorname{rank}(M_a)=2,
\qquad
\ker(M_a)=\operatorname{span}\{k\}.
\]
The same lemma and the nonzero factors above give
\[
\det(M_b)=hrtupz^2\ne0,
\]
and \(\det(M_c)\ne0\) is part of the chart definition.  Thus both nonroot slices have rank three, so the rank-two label is uniquely \(a\).

The two sign assumptions translate exactly into
\[
k_ak_b=(-uz)(z+vy)=-uz(z+vy)<0,
\qquad
k_ak_c=(-uz)(-y)=uzy>0.
\]
Thus the \(b\)-coordinate, and only that nonroot coordinate, has sign opposite to the \(a\)-coordinate.  Any nonzero \(\widetilde k\in\ker(M_a)\) equals \(\lambda k\) for some \(\lambda\ne0\); multiplication by \(\lambda\) either preserves all signs or reverses all signs simultaneously.  The sign comparison therefore identifies \(b\) independently of the kernel representative, and \(c\) is the remaining label.

By the principal algebra lemma,
\[
W_b=-Ktyz\ne0,
\qquad
D=K^2ty^2z^2\ne0.
\]
Division is therefore legitimate, and the remaining identities of that lemma give
\[
\frac{W_c}{W_b}
=\frac{-Ktvyz}{-Ktyz}=v,
\qquad
-\frac{N_aW_b}{D}
=\frac{K^2tuy^2z^2}{K^2ty^2z^2}=u.
\]
The structural matrix \(B_{\pi}(u,v)\) is a mechanism parameter, whereas \(\operatorname{Rec}(C)\) is computed from the observed cubic tensor.  Their equality now follows from the two derived ratios and the definition of the recovery rule; it was not imposed as a definition of the target.

Suppose a second principal-chart representation exists.  The rank pattern is intrinsic to the common tensor, so the unique rank-two slice gives \(a'=a\).  The kernel and its sign comparisons are also intrinsic, so \(b'=b\) and then \(c'=c\).  With these common labels, \(N,W,D\) are functions only of \(C\), and the two ratios above force \(u'=u\) and \(v'=v\).  This proves global uniqueness on the union \(\mathcal C^{\circ}\).  Finally, there are only three slices; computing their fixed-size ranks, one adjugate, two matrix-vector products or contractions, finitely many signs, and two ratios has constant cost in the real-arithmetic model.
\end{proof}

\begin{proposition}[prop:exceptional-cubic-chart-boundaries]\label{prop:exceptional-cubic-chart-boundaries}
Two distinct phenomena occur on the complement of the principal chart.  First, for a fixed known order \(\pi=(a,b,c)\), if \(y=0\) and \(hxrtz\ne0\), the zero \(c\)-column of \(\operatorname{adj}(M_a)\) is a removable chart failure: with
\[
k=(-uz,z,0)^{\mathsf T}
\]
one has
\[
M_bk=tz(0,1,v)^{\mathsf T},
\qquad
v=\frac{(M_bk)_c}{(M_bk)_b},
\qquad
u=-\frac{k_a}{k_b+vk_c},
\]
and \(\operatorname{adj}(M_a)e_b\ne0\).  Second, on the exceptional hypersurface \(r=0\), the fixed-order cubic map has genuine coefficient ambiguity: for every \(\Delta\ne0\), the replacement
\[
u'=u+\Delta,
\qquad
y'=y-x\Delta
\]
with all other displayed cubic parameters fixed leaves \(C_{\pi}^{(3)}\) unchanged while changing \(u\).  The second assertion is a cubic-tensor ambiguity and does not assert equality of covariance matrices or equality in the full published second-and-third-moment fiber.
\end{proposition}
\begin{proof}
Assume first that \(y=0\) and \(hxrtz\ne0\).  Then the algebraic kernel vector becomes
\[
k=(-uz,z,0)^{\mathsf T}.
\]
Because \(G^{\mathsf T}k=R^{\mathsf T}k=0\) and \(E^{\mathsf T}k=z\), the slice identity
\(M_b=hpGG^{\mathsf T}+ruRR^{\mathsf T}+tEE^{\mathsf T}\) gives
\[
M_bk=tE(E^{\mathsf T}k)=tz(0,1,v)^{\mathsf T}.
\]
The denominator \((M_bk)_b=tz\) is nonzero, so
\[
v=\frac{(M_bk)_c}{(M_bk)_b}.
\]
Also \(k_b+vk_c=z\ne0\), and hence
\[
-\frac{k_a}{k_b+vk_c}=-\frac{-uz}{z}=u.
\]
Writing \(K=hxr\), the principal algebra identity remains valid on this boundary and yields
\[
\operatorname{adj}(M_a)e_b=Kk k_b=Kzk\ne0.
\]
Thus the vanishing of the originally selected \(c\)-column is removable for the fixed-order coefficient calculation.

For the second assertion, set \(r=0\), choose \(\Delta\ne0\), and define \(u'=u+\Delta\) and \(y'=y-x\Delta\).  Then
\[
p'=u'x+y'=(u+\Delta)x+y-x\Delta=ux+y=p
\]
and, because \(v,z\) are fixed,
\[
q'=vp'+z=vp+z=q.
\]
Consequently the global-source vector \((x,p,q)^{\mathsf T}\), the middle-source vector \((0,1,v)^{\mathsf T}\), and the terminal-source vector \((0,0,1)^{\mathsf T}\) are unchanged.  The only remaining vector depending on \(u\) multiplies \(r\), and its tensor contribution vanishes because \(r=0\).  Therefore
\[
C_{\pi}^{(3)}(h,0,t,w,x,y,z,u,v)
=C_{\pi}^{(3)}(h,0,t,w,x,y',z,u',v),
\]
while \(u'\ne u\).  This proves genuine noninjectivity of the fixed-order cubic map on that hypersurface, with the stated limitation concerning second moments.
\end{proof}

\begin{proposition}[prop:analytic-full-ideal-membership]\label{prop:analytic-full-ideal-membership}
For the exact published moment parameterization, the polynomial \(R_{\mathrm{el}}=\beta W_2^{\mathrm{el}}-W_3^{\mathrm{el}}\) vanishes identically after substitution of the moment formulas and belongs to \(\mathcal I_{\leq3}\). Consequently, Gröbner reduction of \(R_{\mathrm{el}}\) by any genuine Gröbner basis of \(\mathcal I_{\leq3}\), for any fixed monomial order, has zero remainder. This is an analytic ideal-membership conclusion and does not assert that the released computation was independently executed.
\end{proposition}
\begin{proof}
Write \(h_j=\omega_j^{(3)}\) for \(0\leq j\leq3\), and introduce the proof intermediates
\[
p=u\gamma _1+\gamma _2,\qquad q=\beta p+\gamma _3,
\]
\[
g=(\gamma _1,p,q)^{\mathsf T},\qquad
d=(1,u,\beta u)^{\mathsf T},\qquad
e=(0,1,\beta)^{\mathsf T},\qquad
f=(0,0,1)^{\mathsf T}.
\]
By the displayed matrix in Definition~\(\mathrm{def{:}published\text{-}mixing\text{-}map}\), its four columns are \(g,d,e,f\).  Hence the cubic moment equations in that definition give the polynomial identity
\[
C=h_0g^{\otimes3}+h_1d^{\otimes3}+h_2e^{\otimes3}+h_3f^{\otimes3}.
\]
In particular, the first and second cubic slices are
\[
M_1=h_0\gamma _1gg^{\mathsf T}+h_1dd^{\mathsf T},
\qquad
M_2=h_0pgg^{\mathsf T}+h_1u,dd^{\mathsf T}+h_2ee^{\mathsf T}.
\]
The three entries in Definition~\(\mathrm{def{:}linear\text{-}elimination\text{-}certificate}\) are precisely the cofactors in the third column of \(\operatorname{adj}(M_1)\).  Expanding those three two-by-two determinants, with no division and no nonvanishing hypothesis, gives
\[
N^{\mathrm{el}}
=h_0h_1\gamma _1
\begin{pmatrix}
u\gamma _2\gamma _3\\
-\gamma _2(\gamma _3+\beta\gamma _2)\\
\gamma _2^2
\end{pmatrix}.
\]
Set \(K=h_0h_1\gamma _1\).  Using only \(p=u\gamma _1+\gamma _2\) and \(q=\beta p+\gamma _3\), direct contraction yields
\[
\begin{aligned}
g^{\mathsf T}N^{\mathrm{el}}
&=K\{u\gamma _1\gamma _2\gamma _3
-p\gamma _2(\gamma _3+\beta\gamma _2)+q\gamma _2^2\}\\
&=K\gamma _2\gamma _3(u\gamma _1-p+\gamma _2)=0,\\
d^{\mathsf T}N^{\mathrm{el}}
&=Ku\{\gamma _2\gamma _3-\gamma _2(\gamma _3+\beta\gamma _2)
+\beta\gamma _2^2\}=0,\\
e^{\mathsf T}N^{\mathrm{el}}
&=K\{-\gamma _2(\gamma _3+\beta\gamma _2)+\beta\gamma _2^2\}
=-K\gamma _2\gamma _3.
\end{aligned}
\]
Substitution in the formula for \(M_2\) therefore gives, equation by equation,
\[
\begin{aligned}
M_2N^{\mathrm{el}}
&=h_0p,g(g^{\mathsf T}N^{\mathrm{el}})
+h_1u,d(d^{\mathsf T}N^{\mathrm{el}})
+h_2e(e^{\mathsf T}N^{\mathrm{el}})\\
&=-h_0h_1h_2\gamma _1\gamma _2\gamma _3(0,1,\beta)^{\mathsf T}.
\end{aligned}
\]
By the definitions of \(W_2^{\mathrm{el}}\) and \(W_3^{\mathrm{el}}\), these are respectively the second and third coordinates of \(M_2N^{\mathrm{el}}\).  Thus
\[
\beta W_2^{\mathrm{el}}-W_3^{\mathrm{el}}=0
\]
identically under the exact published parameterization.

It remains to turn this parameter identity into ideal membership.  Let
\[
T=\mathbb F[\beta,\gamma _1,\gamma _2,\gamma _3,u,
\omega_0^{(2)},\ldots,\omega_3^{(2)},
\omega_0^{(3)},\ldots,\omega_3^{(3)}]
\]
and let \(\rho:\mathcal R_{\mathrm{pub}}\to T\) fix every parameter variable and replace each of the sixteen moment variables by its formula from Definition~\(\mathrm{def{:}published\text{-}mixing\text{-}map}\).  Every generator of \(\mathcal J_{\leq3}\) maps to zero, so
\[
\mathcal J_{\leq3}\subseteq\ker\rho.
\]
Conversely, replace the moment variables one at a time in any polynomial \(F\in\mathcal R_{\mathrm{pub}}\).  At a replacement of a moment variable \(r\) by its parameter polynomial \(r_0\), the identity
\[
r^k-r_0^k=(r-r_0)\sum_{j=0}^{k-1}r^{k-1-j}r_0^j
\]
shows that the change lies in the ideal generated by \(r-r_0\).  After all sixteen replacements this proves
\[
F-\rho(F)\in\mathcal J_{\leq3}.
\]
If \(F\in\ker\rho\), the right term is zero, hence \(F\in\mathcal J_{\leq3}\).  Therefore
\[
\ker\rho=\mathcal J_{\leq3}.
\]
The identity already proved says \(\rho(R_{\mathrm{el}})=0\).  Since \(R_{\mathrm{el}}\in\mathbb F[\beta,\mathbf m]\), Definitions~\(\mathrm{def{:}full\text{-}published\text{-}moment\text{-}ideal}\) and \(\mathrm{def{:}published\text{-}elimination\text{-}ideal}\) now imply
\[
R_{\mathrm{el}}\in
\mathcal J_{\leq3}\cap\mathbb F[\beta,\mathbf m]
=\mathcal I_{\leq3}.
\]

Finally, let \(\mathcal H\) be any Gröbner basis of \(\mathcal I_{\leq3}\) for a fixed monomial order, and let \(r\) be a remainder obtained by Gröbner reduction of \(R_{\mathrm{el}}\) by \(\mathcal H\).  Then \(r\in\mathcal I_{\leq3}\), because both \(R_{\mathrm{el}}\) and the subtracted combination of elements of \(\mathcal H\) lie in that ideal.  If \(r\neq0\), its leading monomial lies in the initial ideal of \(\mathcal I_{\leq3}\).  By the defining property of a Gröbner basis, that initial ideal is generated by the leading monomials of \(\mathcal H\), so one of them divides the leading monomial of \(r\).  This contradicts the defining property of a reduced remainder.  Hence \(r=0\), without any claim that a particular stored basis was run.
\end{proof}

\section*{Technical internal limitation (diagnostic only)}

The exact released graph, the \(3\times4\) mixing convention, all six covariance coordinates, all ten cubic coordinates, and the proposed linear relation are specified in the formal core. The pinned Macaulay2 run has not yet printed the entire stored basis reduction of \(R_{\mathrm{el}}\) or the localized target degree. Zero remainder against that full basis is the first proof checkpoint: a nonzero remainder or a retained generic second full-system root ends the localized re-audit claim. Until this checkpoint succeeds, the independent symbolic derivation does not supersede the conclusion of Theorem 3.6. The unrestricted real atlas is also not claimed: known zero-skew twins and removable zero-column chart failures show that a denominator divisor alone cannot equal the nonidentification locus.

\section{Honest scope}

The theorem targets concern global recovery on the compact positive chamber, a pinned and explicitly gated full-ideal re-audit on the fixed published graph, an explicit covariance-overlap neighborhood, and exact moment-set coverage. The localized degree-one target stands only if the full-basis zero-remainder checkpoint succeeds. The complete classification of all real fibers outside the chamber remains open. As a sanity reduction, on the fixed order \(1\to2\to3\) the cubic functional returns the same structural \(T\to Y\) coordinate targeted by the published fourth-order functional; this is a direct compatibility corollary, not a separate root-solving result. With no latent factor the structural target reduces to ordinary fully observed LiNGAM recovery as in Shimizu et al. (2006), although the displayed latent-slice chart is then unnecessary.

\begin{thebibliography}{99}

  \bibitem{tramontanoEtAl2025HigherOrderCumulants} D. Tramontano, Y. Kivva, S. Salehkaleybar, M. Drton, and N. Kiyavash (2025), Causal Effect Identification in lvLiNGAM from Higher-Order Cumulants, Proceedings of Machine Learning Research 267, ICML; arXiv:2506.05202.

  \bibitem{caiEtAl2023HigherOrderCumulants} R. Cai, Z. Huang, W. Chen, Z. Hao, and K. Zhang (2023), Causal Discovery with Latent Confounders Based on Higher-Order Cumulants, Proceedings of Machine Learning Research 202, ICML; arXiv:2305.19582.

  \bibitem{garciaPuenteSpielvogelSullivant2010ComputerAlgebra} L. D. Garcia-Puente, S. Spielvogel, and S. Sullivant (2010), Identifying Causal Effects with Computer Algebra, Proceedings of UAI 2010; arXiv:1007.3784.

  \bibitem{drtonFoygelSullivant2011GlobalIdentifiability} M. Drton, R. Foygel, and S. Sullivant (2011), Global Identifiability of Linear Structural Equation Models, Annals of Statistics 39, 865--886; DOI:10.1214/10-AOS859.

  \bibitem{schkodaRobevaDrton2024CausalDiscovery} D. Schkoda, E. Robeva, and M. Drton (2024), Causal Discovery of Linear Non-Gaussian Causal Models with Unobserved Confounding, public manuscript; arXiv:2408.04907.

  \bibitem{tramontanoEtAl2024EffectIdentification} D. Tramontano, Y. Kivva, S. Salehkaleybar, M. Drton, and N. Kiyavash (2024), Causal Effect Identification in LiNGAM Models with Latent Confounders, Proceedings of Machine Learning Research 235, ICML; arXiv:2406.02049.

  \bibitem{ghassamiEtAl2020DistributionEquivalence} A. Ghassami, A. Yang, N. Kiyavash, and K. Zhang (2020), Characterizing Distribution Equivalence and Structure Learning for Cyclic and Acyclic Directed Graphs, Proceedings of Machine Learning Research 119, ICML; arXiv:1910.12993.

  \bibitem{foygelBarberEtAl2022LFHTC} R. Foygel Barber, M. Drton, N. Sturma, and L. Weihs (2022), Half-Trek Criterion for Identifiability of Latent Variable Models, Annals of Statistics 50, 3174--3196; DOI:10.1214/22-AOS2221.

  \bibitem{babii2020HonestConfidenceSets} A. Babii (2020), Honest Confidence Sets in Nonparametric IV Regression and Other Ill-Posed Models, Econometric Theory 36, 658--706; arXiv:1611.03015.

  \bibitem{andrewsGuggenberger2019SingularityRobust} D. W. K. Andrews and P. Guggenberger (2019), Identification- and Singularity-Robust Inference for Moment Condition Models, Quantitative Economics 10, 1703--1746; DOI:10.3982/QE1219.

  \bibitem{anandkumarEtAl2014TensorDecompositions} A. Anandkumar, R. Ge, D. Hsu, S. M. Kakade, and M. Telgarsky (2014), Tensor Decompositions for Learning Latent Variable Models, Journal of Machine Learning Research 15, 2773--2832; arXiv:1210.7559.

  \bibitem{shimizuEtAl2006LiNGAM} S. Shimizu, P. O. Hoyer, A. Hyvarinen, and A. Kerminen (2006), A Linear Non-Gaussian Acyclic Model for Causal Discovery, Journal of Machine Learning Research 7, 2003--2030.

  \bibitem{hoeffding1963ProbabilityInequalities} W. Hoeffding (1963), Probability Inequalities for Sums of Bounded Random Variables, Journal of the American Statistical Association 58, 13--30.

\end{thebibliography}

\end{document}